\providecommand{\toplevelprefix}{../..}  %
\documentclass[../../book-main_zh.tex]{subfiles}

\begin{document}

\chapter{熵、扩散、去噪与有损编码}\label{app:entropy}

\begin{quote}
    ``{\em 随着时间推移，无序度或熵的增加是所谓时间箭头的一个例子，它区分了过去与未来，赋予了时间一个方向。}''

    $~$\hfill ——《时间简史》，斯蒂芬·霍金 (Stephen Hawking)
\end{quote}
\vspace{5mm}

在本附录中，我们将为\Cref{ch:denoising}中提到的几个事实提供证明，这些事实涉及微分熵、它在扩散过程中的演化方式，以及它与有损编码的联系。我们将对表示数据源的随机变量（记为\(\vx\)）作出以下温和的假设。

\begin{assumption}\label{assumption:entropy_x_compact_support}
    \(\vx\)的支撑集是一个半径至多为\(R\)的紧集\(\cS \subseteq \R^{D}\)，即\(R \doteq \sup_{\vxi \in \cS}\norm{\vxi}_{2}\)。
\end{assumption}

特别地，由于欧几里得空间中的紧集是有界的，因此有\(R < \infty\)。我们将始终使用记号\(B_{r}(\vxi) \doteq \{\vu \in \R^{D} \colon \norm{\vxi - \vu}_{2}
\leq r\}\)来表示以\(\vxi\)为中心、半径为\(r\)的欧几里得球。在这个意义上，\Cref{assumption:entropy_x_compact_support}意味着\(\cS \subseteq B_{R}(\vzero)\)。

请注意，这一假设对我们在实践中关心的（几乎）所有变量都成立，因为它（通常）是由数据预处理期间的归一化步骤所施加的。

\section{低维分布的微分熵}\label{sec:low_dim_entropy}

在这个简短的附录中，我们讨论低维分布的微分熵。根据定义，没有密度的随机变量\(\vx\)的微分熵为\(-\infty\)；这包括所有支撑在低维集上的随机变量。本节的目的是讨论为什么这是一个“本质上正确”的值。

事实上，设\(\vx\)为满足\Cref{assumption:entropy_x_compact_support}的任意随机变量，\(\vx\)的支撑集\(\cS\)的体积为\(0\)。\footnote{形式上，这意味着\(\cS\)是Borel可测的，且Borel测度为\(0\)。} 我们将考虑\(\vx\)在\(\cS\)上均匀分布的情况。\footnote{例如，关于\(\cS\)上的Hausdorff测度。} 我们的目标是计算\(h(\vx)\)。

在这种情况下，\(\vx\)将没有密度；在假设我们不知道\(h(\vx) = -\infty\)的反事实世界中，我们无法使用微分熵的标准定义直接定义它。相反，正如在分析学和信息论的其他部分一样，考虑\(\vx\)的具有密度的逐次更优近似\(\vx_{\eps}\)的熵的极限是合理的，即：
\begin{equation}
    h(\vx)\ \text{``=''}\ \lim_{\eps \searrow 0}h(\vx_{\eps}).
\end{equation}

为此，基本思想是取\(\cS\)的\(\eps\)-加厚（\(\eps\)-thickening），记为\(\cS_{\eps}\)，定义为
\begin{equation}
    \cS_{\eps} = \bigcup_{\vxi \in \cS}B_{\eps}(\vxi)
\end{equation}
并在\Cref{fig:entropy_eps_thickening}中进行了可视化。
\begin{figure}[th]
    \centering
    \begin{tikzpicture}
        \def\radius{0.2cm} %
        \def\curve{(0,0) .. controls (1,1.5) and (3,-0.5) .. (4,1)} %

        \path[decoration={markings, mark=between positions 0 and 1
                step 0.02 with {
                    \fill[red!30, opacity=0.5, draw=none] (0,0)
                    circle (\radius);
        }}, postaction={decorate}] \curve;

        \draw[blue, thick] \curve node[right, blue] {$\cS$};

        \coordinate (p_on_curve) at (2, 0.5); %
        \fill[red!50, opacity=0.7] (p_on_curve) circle (\radius);
        \fill[black] (p_on_curve) circle (1pt) node[below left] {$x$};

        \node[red, below right] at (3.5, -0.2) {$\cS_{\eps} =
        \bigcup_{\vxi \in \cS} B_{\eps}(\vxi)$};

    \end{tikzpicture}
    \caption{曲线\(\cS \subseteq \R^{2}\)的\(\eps\)-加厚\(\cS_\eps\)示意图。}
    \label{fig:entropy_eps_thickening}
\end{figure}
我们将处理支撑集为\(\cS_{\eps}\)（全维的）的随机变量，并取\(\eps \to 0\)的极限。实际上，定义\(\vx_{\eps} \sim \dUnif(\cS_{\eps})\)。由于\(\cS_{\eps}\)具有正体积，\(\vx_{\eps}\)具有密度\(p_{\eps}\)，等于
\begin{equation}
    p_{\eps}(\vxi) = \indvar(\vxi \in \cS_{\eps}) \cdot
    \frac{1}{\volume(\cS_{\eps})}.
\end{equation}
使用\(0 \log 0 = 0\)的约定计算\(\vx_{\eps}\)的熵，可得
\begin{align}
    h(\vx_{\eps})
    &= -\int_{\R^{D}}p_{\eps}(\vxi)\log p_{\eps}(\vxi)\odif{\vxi} \\
    &= -\int_{\cS_{\eps}}\frac{1}{\volume(\cS_{\eps})}
    \log\rp{\frac{1}{\volume(\cS_{\eps})}}\odif{\vxi} \\
    &=
    \frac{\log(\volume(\cS_{\eps}))}{\volume(\cS_{\eps})}\int_{\cS_{\eps}}\odif{\vxi}
    \\
    &= \log(\volume(\cS_{\eps})).
\end{align}
由于\(\cS\)是紧集，\(\volume(\cS_{\eps})\)是有限的，并且当\(\eps \searrow 0\)时趋于\(0\)。因此
\begin{equation}
    h(\vx) = \lim_{\eps \searrow 0}h(\vx_{\eps}) = \lim_{\eps
    \searrow 0}\log(\volume(\cS_{\eps})) = -\infty,
\end{equation}
如我们所愿。

上述计算实际上是一组关于在对\(\vx\)的分布施加某些约束时\(\vx\)的最大可能熵的更为著名和广受赞誉的结果的推论。如果我们不在这里提供这些结果，那将是我们的疏忽；例如，\cite{poliyanski2024information}的第2章提供了这些证明。
\begin{theorem}\label{thm:max_entropy}
    设\(\vx\)为\(\R^{D}\)上的随机变量。
    \begin{enumerate}
        \item 如果\(\vx\)的支撑集是紧集\(\cS \subseteq \R^{D}\)（即\Cref{assumption:entropy_x_compact_support}），那么
            \begin{equation}
                h(\vx) \leq h(\dUnif(\cS)) = \log \volume(\cS).
            \end{equation}
        \item 如果\(\vx\)具有有限协方差，使得对于半正定（PSD）矩阵\(\vSigma \in \PSD(D)\)，有\(\Cov(\vx) \preceq \vSigma\)（关于PSD偏序，即\(\vSigma - \Cov(\vx)\)是PSD的），那么
            \begin{equation}
                h(\vx) \leq h(\dNorm(\vzero, \vSigma)) =
                \frac{1}{2}\log((2\pi e)^{D}\det\vSigma).
            \end{equation}
        \item 如果\(\vx\)具有有限二阶矩，使得对于常数\(a \geq 0\)，有\(\Ex\norm{\vx}_{2}^{2} \leq a\)，那么
            \begin{equation}
                h(\vx) \leq h\rp{\dNorm\rp{\vzero, \frac{a}{D}\vI}} =
                \frac{D}{2}\log\frac{2\pi e a}{D}.
            \end{equation}
    \end{enumerate}
\end{theorem}

\section{扩散与去噪过程}\label{sec:entropy_diffusion}

在正文（\Cref{ch:denoising}）中，我们考虑了一个随机变量\(\vx\)，以及由\eqref{eq:additive_gaussian_noise_model}定义的随机过程，即
\begin{equation}\label{eq:app_additive_gaussian_noise_model}
    \vx_{t} = \vx + t\vg, \qquad  \forall t \in [0, T]
\end{equation}
其中\(\vg \sim \dNorm(\vzero, \vI)\)且独立于\(\vx\)。

本节的结构如下。在\Cref{sub:diffusion_entropy_increases}中，我们提供了一个正式的定理和清晰的证明，表明在\Cref{eq:app_additive_gaussian_noise_model}下熵是增加的，即\(\odv*{h(\vx_{t})}{t} > 0\)。在\Cref{sub:denoising_entropy_decreases}中，我们提供了一个正式的定理和清晰的证明，表明在\Cref{eq:app_additive_gaussian_noise_model}下，去噪过程中的熵是减少的，即对于所有\(s < t\)，有\(h(\Ex[\vx_{s} \given \vx_{t}]) < h(\vx_{t})\)。在\Cref{sub:app_diffusion_intermediate_results}中，我们为建立前几个小节中的主张所需的技术引理提供了证明。

在开始之前，我们介绍一些关键记号。首先，设\(\phi_{t}\)为\(\dNorm(\vzero, t^{2}\vI)\)的密度，即
\begin{equation}\label{eq:gaussian_noise_time_t}
    \phi_{t}(\vxi) \doteq
    \frac{1}{(2\pi)^{D/2}t^{D}}\exp\rp{-\frac{\norm{\vxi}_{2}^{2}}{2t^{2}}}.
\end{equation}
其次，\(\vx_{t}\)的支撑集为整个\(\R^{D}\)，因此它具有\textit{密度}，我们将其记为\(p_{t}\)（与正文中一样）。简单的计算表明
\begin{equation}\label{eq:p_t_representation}
    p_{t}(\vxi) = \Ex[\phi_{t}(\vxi - \vx)],
\end{equation}
从这个表示中可以推导出（即根据\Cref{prop:diff_convolution}），\(p_{t}\)关于\(\vxi\)是光滑的（即无限可微），关于\(t\)也是光滑的，并且处处为正。乍一看，这个事实有些引人注目：即使对于一个完全不规则的随机变量\(\vx\)（例如，没有密度的伯努利随机变量），其高斯平滑在每个（任意小的）\(t > 0\)处都具有密度。该证明留给精通数学分析的读者作为练习。

然而，我们还需要增加一个关于\(\vx\)分布的\textit{光滑性}的假设，这将消除在低维分布中\(t = 0\)附近出现的一些技术问题。\footnote{因为那时各种量会变得高度不规则，处理它们将需要大量额外的分析。} 尽管如此，我们预计通过额外的工作，我们的结果在更温和的假设下依然成立。现在，让我们假设：
\begin{assumption}\label{assumption:entropy_x_density}
    \(\vx\)具有二次连续可微的密度，记为\(p\)。
\end{assumption}

\subsection{扩散过程随时间增加熵}\label{sub:diffusion_entropy_increases}

在本节附录中，我们提供\Cref{thm:diffusion_entropy_increases}的证明。为方便起见，我们将其重述如下。

\begin{theorem}[扩散增加熵]\label{thm:diffusion_entropy_increases}
    设\(\vx\)为满足\Cref{assumption:entropy_x_compact_support,assumption:entropy_x_density}的任意随机变量，并设\((\vx_{t})_{t \in [0, T]}\)为随机过程\eqref{eq:app_additive_gaussian_noise_model}。那么
    \begin{equation}\label{eq:diffusion_entropy_increases}
        h(\vx_{s}) < h(\vx_{t}), \qquad \forall s, t \colon 0 \leq s < t \leq T.
    \end{equation}
\end{theorem}
\begin{proof}
    在开始之前，我们必须问：\eqref{eq:diffusion_entropy_increases}中的不等式何时有意义？我们将在\Cref{lem:diffusion_entropy_exists}中证明，在我们的假设下，微分熵是良定义的，永远不会是\(+\infty\)，并且对于\(t > 0\)是有限的，因此\eqref{eq:diffusion_entropy_increases}中的（严格）不等式是有意义的。

    撇开良定义的问题不谈，本证明的关键是证明\(\vx_{t}\)的密度\(p_{t}\)满足一个特定的偏微分方程，该方程与\textit{热传导方程}非常相似。热传导方程是一个著名的偏微分方程（PDE），描述了热量在空间中的扩散。这在直觉上应该是合理的，并描绘了一幅心理图像：随着时间\(t\)的增加，来自原始（可能高度集中）\(\vx\)的概率像真空中热源辐射的热量一样分散到整个\(\R^{D}\)中。

    对于更一般的随机过程，这种关于\(p_{t}\)的偏微分方程被称为\textit{福克-普朗克方程（Fokker-Planck equations）}，它们是非常强大的工具，因为它们允许我们用\(p_{t}\)的瞬时空间导数来描述\(p_{t}\)的瞬时时间导数，反之亦然，从而提供了对\(p_{t}\)的正则性和动力学的简洁描述。一旦我们获得了\(p_{t}\)的动力学，我们就可以利用该系统获得另一个描述\(h(\vx_{t})\)动力学的系统，毕竟\(h(\vx_{t})\)只是\(p_{t}\)的一个泛函。

    该偏微分方程的描述涉及一个称为拉普拉斯算子（Laplacian）\(\Delta\)的数学对象。回想一下多元微积分课程，作用于时间上可微且空间上二次可微的函数\(f \colon (0, T) \times \R^{D} \to \R\)的拉普拉斯算子由下式给出：
    \begin{equation}
        \Delta f_{t}(\vxi) = \tr(\nabla^{2}f_{t}(\vxi)) = \sum_{i =
        1}^{D}\pddv{f_{t}}{\xi_{i}}(\vxi).
    \end{equation}

    也就是说，通过使用\(p_{t}\)的积分表示并在积分号下求导，我们可以计算\(p_{t}\)的导数（我们在\Cref{prop:p_t_derivatives}中进行了计算），并观察到\(p_{t}\)满足类似热传导的偏微分方程：
    \begin{equation}
        \pdv{p_{t}}{t}(\vxi) = t\Delta p_{t}(\vxi).
    \end{equation}
    然后，为了寻找\(h(\vx_{t})\)的动力学，我们可以再次使用\Cref{prop:dutis}以及类似热传导的偏微分方程，得到
    \begin{align}
        \odv*{h(\vx_{t})}{t}
        &= -\odv*{\int_{\R^{D}}p_{t}(\vxi)\log p_{t}(\vxi)\odif{\vxi}}{t} \\
        &= -\int_{\R^{D}}\pdv*{\bs{p_{t}(\vxi)\log
        p_{t}(\vxi)}}{t}\odif{\vxi} \\
        &= -\int_{\R^{D}}\pdv{p_{t}}{t}(\vxi)[1 + \log
        p_{t}(\vxi)]\odif{\vxi} \\
        &= -t\int_{\R^{D}}\Delta p_{t}(\vxi)[1 + \log p_{t}(\vxi)]\odif{\vxi}.
    \end{align}
    通过使用一个稍微复杂的分部积分论证（\Cref{lem:diffusion_ibp}），我们得到
    \begin{align}
        \odv*{h(\vx_{t})}{t}
        &= t\int_{\R^{D}}\ip{\nabla \log p_{t}(\vxi)}{\nabla
        p_{t}(\vxi)}\odif{\vxi} \\
        &= t\int_{\R^{D}}\frac{\norm{\nabla
        p_{t}(\vxi)}_{2}^{2}}{p_{t}(\vxi)}\odif{\vxi} \\
        &> 0
    \end{align}
    其中最后一行严格不等式成立，因为如果它不成立，\(\nabla p_{t}(\vxi)\)将需要几乎处处为零（即除了可能在一个零体积集上之外处处为零），但这将意味着\(p_{t}\)几乎处处为常数，这与\(p_{t}\)是密度的事实相矛盾。

    为了完成证明，我们只需使用微积分基本定理
    \begin{equation}
        h(\vx_{t}) = h(\vx_{s}) +
        \int_{s}^{t}\odv*{h(\vx_{u})}{u}\odif{u} > h(\vx_{s}),
    \end{equation}
    这就证明了该主张。（注意，当\(h(\vx_{s}) = -\infty\)时这没有意义，这只能在\(s = 0\)且\(h(\vx) = -\infty\)时发生，但在这种情况下\(h(\vx_{t}) > -\infty\)，因此该主张无论如何都是空真（vacuously true）的。）
\end{proof}

\subsection{去噪过程随时间减少熵}\label{sub:denoising_entropy_decreases}

回想一下，在\Cref{sub:intro_diffusion_denoising}中，我们从随机变量\(\vx_{T}\)开始，并使用以下形式的迭代对其进行迭代去噪：
\begin{equation}\label{eq:app_denoising_iteration}
    \hat{\vx}_{s} \doteq \Ex[\vx_{s} \mid \vx_{t} = \hat{\vx}_{t}] =
    \frac{s}{t}\hat{\vx}_{t} + \bp{1 -
    \frac{s}{t}}\bar{\vx}^{\ast}(t, \hat{\vx}_{t}).
\end{equation}
其中\(s, t \in \{t_{0}, t_{1}, \dots, t_{L}\}\)且\(s < t\)，\(\vx_{T} = \hat{\vx}_{T}\)。我们希望证明\(h(\hat{\vx}_{s}) < h(\hat{\vx}_{t})\)，从而表明去噪过程实际上减少了熵。

在着手进行证明之前，我们对问题陈述作几点说明。首先，Tweedie公式\eqref{eq:tweedie}表明
\begin{equation}
    \bar{\vx}^{\ast}(t, \vx_{t}) = \vx_{t} + t^{2}\nabla \log p_{t}(\vx_{t}),
\end{equation}
它将从时间\(t\)到时间\(0\)的完整去噪步骤比作在\(\vx_{t}\)的对数密度上的梯度步。我们能否在\eqref{eq:app_denoising_iteration}中得到从时间\(t\)到时间\(s\)的完整去噪步骤的类似结果？事实证明我们确实可以，而且非常简单。通过使用\eqref{eq:app_denoising_iteration}和Tweedie公式\eqref{eq:tweedie}，我们得到
\begin{equation}\label{eq:app_denoising_iteration_score}
    \Ex[\vx_{s} \mid \vx_{t}] = \frac{s}{t}\vx_{t} + \bp{1 -
    \frac{s}{t}}\bp{\vx_{t} + t^{2}\nabla_{\vx_{t}}\log
    p_{t}(\vx_{t})} = \vx_{t} + \bp{1 -
    \frac{s}{t}}t^{2}\nabla_{\vx_{t}}\log p_{t}(\vx_{t}).
\end{equation}
因此，这个迭代去噪步骤再次是在扰动对数密度\(\log p_{t}\)上的梯度步，只是步长缩小了。特别地，这一步可以看作是\textit{得分函数向量场}对随机变量\(\vx_{t}\)分布的扰动，这暗示了与随机微分方程（SDE）和扩散模型理论\cite{song2020score}的联系。事实上，可以使用这种强大的机制和极限论证（例如，遵循\cite{DBLP:conf/iclr/ChenC0LSZ23}阐述中的技术方法）来推导以下结果\Cref{thm:conditioning_reduces_entropy}的证明。我们将在这里给出一个更简单的证明，它只使用初等工具，从而阐明通过去噪减少熵的过程背后的一些关键量。另一方面，由于\Cref{thm:conditioning_reduces_entropy}中的去噪过程并\textit{不}完全对应于生成观测值\(\vx_{t}\)的加噪过程相关的\textit{逆向}过程，我们将需要处理一些稍微技术性的计算。\footnote{对于熟悉扩散模型的读者，我们在这里指的是时间反转的前向过程与由\Cref{thm:conditioning_reduces_entropy}定义的过程生成的迭代序列不一致。在某种极限下，即采取无限多步的\Cref{thm:conditioning_reduces_entropy}且每步添加无限小水平的噪声时，这些过程是一致的；对于一般的有限步，无论我们的工具多么复杂，我们都必须引入一些近似。}

我们想要证明\(h(\Ex[\vx_{s} \mid \vx_{t}]) < h(\vx_{t})\)，即形式上：
\begin{theorem}\label{thm:conditioning_reduces_entropy}
    设\(\vx\)为满足\Cref{assumption:entropy_x_compact_support,assumption:entropy_x_density}的任意随机变量，并设\((\vx_{t})_{t \in [0, T]}\)为随机过程\eqref{eq:app_additive_gaussian_noise_model}。那么
    \begin{equation}
        h(\Ex[\vx_{s} \mid \vx_{t}]) < h(\vx_{t}),
    \end{equation}
 对于所有$s,  t \in [0, T]$：
\begin{equation}
 \quad 0 < t \leq \frac{R}{\sqrt{2D}},
        \quad 0 \leq s  < t\cdot\min\bc{1, \frac{R^{2}/D -
        2t^{2}}{R^{2}/D - t^{2}}}.
    \end{equation}
\end{theorem}
\begin{proof}
    该证明使用了两个主要思想：
    \begin{enumerate}
        \item 首先，使用变量代换公式写出\(\Ex[\vx_{s} \mid \vx_{t}]\)的密度。
        \item 其次，界定该密度以控制熵。
    \end{enumerate}
    变量代换由\Cref{cor:gribonval_A2}证明其合理性，该推论最初在\cite{Gribonval2011-pf}中推导得出。

    我们现在执行这些想法。由\Cref{cor:gribonval_A2}，我们得到定义为\(\bar{\vx}(\vxi) \doteq \Ex[\vx_{s} \given \vx_{t} = \vxi]\)的函数\(\bar{\vx}\)是可微的、单射的，因此在其值域（我们此后将其记为\(\cX \subseteq \R^{D}\)）上是可逆的。我们将其逆记为\(\bar{\vx}^{-1}\)。使用变量代换公式，\(\bar{\vx}(\vx_{t})\)的密度\(\bar{p}\)由下式给出：
    \begin{equation}
        \bar{p}(\vxi) \doteq \frac{(p_{t} \circ
        \bar{\vx}^{-1})(\vxi)}{\det(\bar{\vx}^{\prime}(\bar{\vx}^{-1}(\vxi)))},
    \end{equation}
    其中（回想一下\Cref{sec:autodiff}）\(\bar{\vx}^{\prime}\)是\(\bar{\vx}\)的雅可比矩阵。由于从\Cref{lem:gribonval_A1}我们知道\(\bar{\vx}^{\prime}\)是正定矩阵，其行列式为正，因此整个密度为正。由此可得
    \begin{align}
        h(\bar{\vx}(\vx_{t}))
        &= -\int_{\cX}\frac{(p_{t} \circ
        \bar{\vx}^{-1})(\vxi)}{\det(\bar{\vx}^{\prime}(\bar{\vx}^{-1}(\vxi)))}
        \log \frac{(p_{t} \circ
        \bar{\vx}^{-1})(\vxi)}{\det(\bar{\vx}^{\prime}(\bar{\vx}^{-1}(\vxi)))}
        \odif{\vxi} \\
        &= -\int_{\cX}\frac{(p_{t} \circ
        \bar{\vx}^{-1})(\vxi)}{\det(\bar{\vx}^{\prime}(\bar{\vx}^{-1}(\vxi)))}
        \log((p_{t} \circ \bar{\vx}^{-1})(\vxi))\odif{\vxi} \\
        &\qquad + \int_{\cX}\frac{(p_{t} \circ
        \bar{\vx}^{-1})(\vxi)}{\det(\bar{\vx}^{\prime}(\bar{\vx}^{-1}(\vxi)))}\log\det\rp{\bar{\vx}^{\prime}(\bar{\vx}^{-1}(\vxi))}\odif{\vxi}
        \\
        &= -\int_{\R^{D}}p_{t}(\vxi)\log p_{t}(\vxi)\odif{\vxi}
        + \int_{\cX}\frac{(p_{t} \circ
        \bar{\vx}^{-1})(\vxi)}{\det(\bar{\vx}^{\prime}(\bar{\vx}^{-1}(\vxi)))}\log\det\rp{\bar{\vx}^{\prime}(\bar{\vx}^{-1}(\vxi))}\odif{\vxi}
        \\
        &= h(\vx_{t}) - \int_{\cX}\frac{(p_{t} \circ
        \bar{\vx}^{-1})(\vxi)}{\det(\bar{\vx}^{\prime}(\bar{\vx}^{-1}(\vxi)))}\log\rp{\frac{1}{\det(\bar{\vx}^{\prime}(\bar{\vx}^{-1}(\vxi)))}}\odif{\vxi}.
    \end{align}
    我们将研究最后一项（包括\(-\)号），并证明它是负的。

    由凹性可知，对于每个\(x \geq 0\)，有\(-x\log x \leq 1 - x\)。因此
    \begin{align}
        h(\bar{\vx}(\vx_{t})) - h(\vx_{t})
        &= - \int_{\cX}\frac{(p_{t} \circ
        \bar{\vx}^{-1})(\vxi)}{\det(\bar{\vx}^{\prime}(\bar{\vx}^{-1}(\vxi)))}\log\rp{\frac{1}{\det(\bar{\vx}^{\prime}(\bar{\vx}^{-1}(\vxi)))}}\odif{\vxi}
        \\
        &\leq  \int_{\cX}(p_{t} \circ \bar{\vx}^{-1})(\vxi)\cdot
        \bp{1 -
        \frac{1}{\det(\bar{\vx}^{\prime}(\bar{\vx}^{-1}(\vxi)))}}\odif{\vxi} \\
        &= \int_{\cX}(p_{t} \circ \bar{\vx}^{-1})(\vxi)\odif{\vxi} -
        \int_{\cX}\frac{(p_{t} \circ
        \bar{\vx}^{-1})(\vxi)}{\det(\bar{\vx}^{\prime}(\bar{\vx}^{-1}(\vxi)))}\odif{\vxi}
        \\
        &=
        \int_{\R^{D}}p_{t}(\vxi)\det\rp{\bar{\vx}^{\prime}(\bar{\vx}^{-1}(\vxi))}\odif{\vxi}
        - \int_{\cX}\bar{p}(\vxi)\odif{\vxi} \\
        &= \int_{\R^{D}}p_{t}(\vxi)\det\rp{\vI + \bp{1 -
        \frac{s}{t}}t^{2}\nabla^{2}\log p_{t}(\vxi)}\odif{\vxi} - 1.
    \end{align}
    现在，通过对特征值使用均值不等式（AM-GM不等式），对于任意对称正定矩阵\(\vM \in \PSD(D)\)，我们有以下界限：
    \begin{equation}
        \det(\vM)^{1/D} = \prod_{i = 1}^{D}\lambda_{i}(\vM)^{1/D}
        \leq \frac{\sum_{i = 1}^{D}\lambda_{i}(\vM)}{D} = \frac{\tr(\vM)}{D},
    \end{equation}
    我们可以将其应用于上述表达式并得到
    \begin{align}
        &\int_{\R^{D}}p_{t}(\vxi)\det\rp{\vI + \bp{1 -
        \frac{s}{t}}t^{2}\nabla^{2}\log p_{t}(\vxi)}\odif{\vxi} \\
        &\leq \int_{\R^{D}} p_{t}(\vxi) \tr\rp{\frac{1}{D}\bs{\vI +
        \bp{1 - \frac{s}{t}}t^{2}\nabla^{2}\log p_{t}(\vxi)}}^{D}\odif{\vxi} \\
        &= \int_{\R^{D}} p_{t}(\vxi)\bp{1 + \frac{\bp{1 -
        \frac{s}{t}}t^{2}}{D}\tr(\nabla^{2}\log p_{t}(\vxi))}^{D}\odif{\vxi} \\
        &= \int_{\R^{D}} p_{t}(\vxi)\bp{1 + \frac{\bp{1 -
        \frac{s}{t}}t^{2}}{D}\Delta \log p_{t}(\vxi)}^{D}\odif{\vxi}.
    \end{align}
    由\Cref{lem:app_diffusion_laplacian_control}可知（回想一下，\(R\)是如\Cref{assumption:entropy_x_compact_support}中\(\vx\)支撑集的半径）
    \begin{equation}
        \abs{\Delta \log p_{t}(\vxi)} \leq \max\rp{\frac{D}{t^{2}},
        \abs*{\frac{R^{2}}{t^{4}} - \frac{D}{t^{2}}}} =: U_{t}.
    \end{equation}
    那么有
    \begin{equation}
        -\frac{\bp{1 - \frac{s}{t}}t^{2}}{D}U_{t} \leq \frac{\bp{1 -
        \frac{s}{t}}t^{2}}{D}\Delta \log p_{t}(\vxi) \leq \frac{\bp{1
        - \frac{s}{t}}t^{2}}{D}U_{t}.
    \end{equation}
    同时，函数\(x \mapsto (1 + x)^{D}\)在\([-1, \infty)\)上是凸的，因此对于\(-(1-s/t)t^{2}U_{t}/D \leq x \leq (1-s/t)t^{2}U_{t}/D\)，我们有
    \begin{align}
        (1 + x)^{D}
        &\leq \bp{1 - \frac{\bp{1 - \frac{s}{t}}t^{2}U_{t}}{D}}^{D} +
        \underbrace{\bs{\bp{1 + \frac{\bp{1 -
                \frac{s}{t}}t^{2}U_{t}}{D}}^{D} - \bp{1 - \frac{\bp{1 -
        \frac{s}{t}}t^{2}U_{t}}{D}}^{D}}}_{M(s, t, D)}x \\
        &\leq 1 + M(s, t, D)x.
    \end{align}
    这里\(M(s, t, D) > 0\)，因为\(U_{t} > 0\)。在上述界限中，我们需要验证\(x\)的下界是\(\geq -1\)的。实际上，
    \begin{align}
        -\frac{\bp{1 - \frac{s}{t}}t^{2}}{D}U_{t}
        &= -\frac{\bp{1 -
        \frac{s}{t}}t^{2}}{D}\max\rp{\frac{D}{t^{2}},
        \abs*{\frac{R^{2}}{t^{4}} - \frac{D}{t^{2}}}} \\
        &= -\bp{1 - \frac{s}{t}}\max\rp{1, \abs*{\frac{R^{2}}{Dt^{2}} - 1}}
    \end{align}
    注意，当且仅当\(\bp{1 - \frac{s}{t}}\cdot\bp{\frac{R^{2}}{Dt^{2}} - 1} \geq 1\)时，这\(\geq -1\)，即\(0 < t < R/\sqrt{2D}\)且\(0 \leq s \leq t\cdot\frac{R^{2}/D - 2t^{2}}{R^{2}/D - t^{2}}\)，这由假设保证。

    应用这个界限，我们得到
    \begin{align}
        &\int_{\R^{D}} p_{t}(\vxi)\bp{1 + \frac{\bp{1 -
        \frac{s}{t}}t^{2}}{D}\Delta \log p_{t}(\vxi)}^{D}\odif{\vxi} \\
        &\leq \int_{\R^{D}}p_{t}(\vxi)\bp{1 + M(s, t, D)\Delta \log
        p_{t}(\vxi)}\odif{\vxi} \\
        &= 1 + M(s, t, D)\int_{\R^{D}}p_{t}(\vxi)\Delta \log
        p_{t}(\vxi)\odif{\vxi} \\
        &= 1 - M(s, t, D)\int_{\R^{D}}\ip{\nabla
        p_{t}(\vxi)}{\nabla\log p_{t}(\vxi)}\odif{\vxi} \\
        &= 1 - M(s, t, D)\int_{\R^{D}}\frac{\norm{\nabla
        p_{t}(\vxi)}_{2}^{2}}{p_{t}(\vxi)}\odif{\vxi},
    \end{align}
    其中最后几行与\Cref{thm:diffusion_entropy_increases}的证明相同。将此结果与我们之前的估计相结合，
    \begin{equation}
        h(\bar{\vx}(\vx_{t})) - h(\vx_{t}) \leq - M(s, t,
        D)\int_{\R^{D}}\frac{\norm{\nabla
        p_{t}(\vxi)}_{2}^{2}}{p_{t}(\vxi)}\odif{\vxi} < 0
    \end{equation}
    其中不等式是严格的，理由与\Cref{thm:diffusion_entropy_increases}中的论证相同。
\end{proof}

请注意，\(s\)和\(t\)的界限取决于数据分布的半径\(R\)，并且不如\Cref{thm:diffusion_entropy_increases}中的界限那么普遍。在以下意义上，该结果实际上是“足够普遍的（as general as needed）”。注意，如果\(\vx\)具有支撑在以\(\vzero\)为中心、半径为\(R\)的球上的二次连续可微密度，那么它对于\(2R\)、\(3R\)等也是如此，即对于任何半径为\(R^{\prime} > R\)的球也是如此。因此，获得适当去噪保证的一种策略是：固定数据维度\(D\)和离散化调度，然后（在分析中）将数据半径\(R\)设置得非常大，使得每个去噪步骤都满足\Cref{thm:conditioning_reduces_entropy}中给出的熵减少要求。那么去噪过程的每一步都将确实如预期那样减少熵。

\subsection{技术引理与中间结果}\label{sub:app_diffusion_intermediate_results}

在本小节中，我们提出支持我们两个主要概念性定理的技术结果。我们的陈述将或多或少符合数学的标准；我们将首先从较高层次的结果开始，然后逐渐回到更基础的结果。较高层次的结果将使用基础结果，通过这种方式，我们得到了一个易于阅读的结果依赖顺序：没有结果依赖于它前面的结果。互不依赖的结果通常按照它们在上述两个证明中出现的位置进行排序。

\subsubsection{微分熵的有限性}

我们首先证明熵在随机过程中存在且是有限的。

\begin{lemma}\label{lem:diffusion_entropy_exists}
    设\(\vx\)为任意随机变量，并设\((\vx_{t})_{t \in [0, T]}\)为随机过程\eqref{eq:app_additive_gaussian_noise_model}。
    \begin{enumerate}
        \item 对于\(t > 0\)，微分熵\(h(\vx_{t})\)存在且\(> -\infty\)。
        \item 如果此外\Cref{assumption:entropy_x_compact_support}对\(\vx\)成立，那么\(h(\vx) < \infty\)且\(h(\vx_{t})\ < \infty\)。
    \end{enumerate}
\end{lemma}
\begin{proof}
    为了证明\Cref{lem:diffusion_entropy_exists}.1，我们使用一个经典但繁琐的分析论证。由于\(\vx_{t}\)具有密度，我们可以写出
    \begin{equation}
        h(\vx_{t}) = -\int_{\R^{D}}p_{t}(\vxi)\log p_{t}(\vxi) \odif{\vxi}.
    \end{equation}
    相应地，设\(g \colon \R^{D} \to \R\)定义为
    \begin{equation}
        g(\vxi) \doteq -p_{t}(\vxi)\log p_{t}(\vxi) \implies
        h(\vx_{t}) = \int_{\R^{D}}g(\vxi)\odif{\vxi}.
    \end{equation}
    像分析中界定积分值的通常做法一样，定义
    \begin{eqnarray}
        &&g_{+}(\vxi) \doteq \max(g(\vxi), 0), \quad g_{-}(\vxi) \doteq
        \max(-g(\vxi), 0) \\ &&\implies \quad g = g_{+} - g_{-}\quad
        \text{and} \quad g_{+}, g_{-} \geq 0.
    \end{eqnarray}
    那么
    \begin{equation}
        h(\vx_{t}) = \int_{\R^{D}}g_{+}(\vxi)\odif{\vxi} -
        \int_{\R^{D}}g_{-}(\vxi)\odif{\vxi},
    \end{equation}
    并且这两个积分都保证是非负的，因为它们的被积函数是非负的。

    为了证明\(h(\vx_{t})\)是良定义的，我们需要证明\(\int_{\R^{D}}g_{+}(\vxi)\odif{\vxi} < \infty\)或\(\int_{\R^{D}}g_{-}(\vxi) \odif{\vxi} < \infty\)。为了证明\(h(\vx_{t}) > -\infty\)，只需证明\(\int_{\R^{D}}g_{-}(\vxi)\odif{\vxi} < \infty\)即可。为了界定\(g_{-}\)的积分，我们需要理解量\(g_{-}\)，即我们想要刻画\(g\)何时为负。
    \begin{equation}
        g(\vxi) \leq 0 \iff p_{t}(\vxi)\log p_{t}(\vxi) \geq 0 \iff
        \log p_{t}(\vxi) \geq 0 \iff p_{t}(\vxi) \geq 1.
    \end{equation}
    因此，有
    \begin{equation}
        g_{-}(\vxi) = \indvar(p_{t}(\vxi) \geq 1)\cdot (-g(\vxi)) =
        \indvar(p_{t}(\vxi) \geq 1)p_{t}(\vxi)\log p_{t}(\vxi).
    \end{equation}
    为了界定\(g_{-}(\vxi)\)的积分，我们需要证明\(p_{t}\)“不是太集中”，即\(p_{t}\)不是太大。为了证明这一点，在这种情况下，我们很幸运能够界定函数\(g_{-}(\vxi)\)本身。即，注意到
    \begin{equation}
        \max_{\vxi \in \R^{D}}\phi_{t}(\vxi - \vx) = \phi_{t}(\vzero)
        = \frac{1}{(2\pi)^{D/2}t^{D}} =: C_{t}.
    \end{equation}
    当\(t \to 0\)时它会趋于无穷大，但对于所有有限的\(t\)它是有限的。因此
    \begin{equation}
        p_{t}(\vxi) = \Ex \phi_{t}(\vxi - \vx) \leq \Ex C_{t} = C_{t}.
    \end{equation}
    现在有两种情况。
    \begin{itemize}
        \item 如果\(C_{t} < 1\)，那么\(p_{t}(\vxi) < 1\)，因此指示函数永远不会是\(1\)，因此\(g_{-} \equiv 0\)，其积分也为\(0\)。
        \item 如果\(C_{t} \geq 1\)，那么\(\log C_{t} \geq 0\)，因此由于对数是单调递增的，
            \begin{align}
                \int_{\R^{D}}g_{-}(\vxi)\odif{\vxi}
                &= \int_{\R^{D}}\indvar(p_{t}(\vxi) \geq
                1)p_{t}(\vxi)\log p_{t}(\vxi)\odif{\vxi}  \\
                &= \Ex[\indvar(p_{t}(\vx_{t}) \geq 1)\log p_{t}(\vx_{t})]  \\
                &\leq \Ex[\indvar(p_{t}(\vx_{t}) \geq 1) \log C_{t}] \\
                &= \Pr[p_{t}(\vx_{t}) \geq 1]\log C_{t}.
            \end{align}
    \end{itemize}
    因此我们有\(\int_{\R^{D}}g_{-}(\vxi)\odif{\vxi} < \infty\)，所以微分熵\(h(\vx_{t})\)存在且\(> -\infty\)。

    为了证明\Cref{lem:diffusion_entropy_exists}.2，假设\Cref{assumption:entropy_x_compact_support}成立。我们想要证明\(h(\vx) < \infty\)且\(h(\vx_{t}) < \infty\)。这样做的机制是相同的，并且涉及最大熵结果\Cref{thm:max_entropy}。即，由于\(\vx\)是绝对有界的，它具有有限的协方差，我们将其记为\(\vSigma\)。那么\(\vx_{t}\)的协方差为\(\vSigma + t^{2}\vI\)。因此，\(\vx\)和\(\vx_{t}\)的熵可以被具有各自协方差的正态分布的熵所上界，即\(\log[(2\pi e)^{D}\det(\vSigma)]\)或\(\log[(2\pi e)^{D}\det(\vSigma + t^{2}\vI)]\)，并且两者都\(< \infty\)。
\end{proof}

\subsubsection{De Bruijn恒等式中的分部积分}

最后，我们补充在\Cref{thm:diffusion_entropy_increases,thm:conditioning_reduces_entropy}的证明中提到的分部积分论证。该论证在概念上非常简单，但需要一些技术估计来证明分部积分中的边界项消失。

\begin{lemma}\label{lem:diffusion_ibp}
    设\(\vx\)为满足\Cref{assumption:entropy_x_compact_support,assumption:entropy_x_density}的任意随机变量，并设\((\vx_{t})_{t \in [0, T]}\)为随机过程\eqref{eq:app_additive_gaussian_noise_model}。对于\(t \geq 0\)，设\(p_{t}\)为\(\vx_{t}\)的密度。那么对于常数\(c \in \R\)，有
    \begin{equation}
        \int_{\R^{D}}\Delta p_{t}(\vxi)[c + \log
        p_{t}(\vxi)]\odif{\vxi} = -\int_{\R^{D}}\ip{\nabla \log
        p_{t}(\vxi)}{\nabla p_{t}(\vxi)}\odif{\vxi}.
    \end{equation}
\end{lemma}
\begin{proof}
    该证明的基本思想分为两步：
    \begin{itemize}
        \item 首先，应用格林公式（Green's theorem）在紧集上进行分部积分；
        \item 其次，将该紧集的半径趋于\(+\infty\)，以得到在整个\(\R^{D}\)上的积分。
    \end{itemize}
    格林公式表明，对于任意紧集\(\cK \subseteq \R^{D}\)、二次连续可微函数\(\plainphi \colon \R^{D} \to \R\)以及连续可微函数\(\psi \colon \R^{D} \to \R\)，
    \begin{equation}
        \int_{\cK}\bc{\psi(\vxi) \Delta \plainphi(\vxi) + \ip{\nabla
        \psi(\vxi)}{\nabla \plainphi(\vxi)}}\odif{\vxi} =
        \int_{\partial \cK}\psi(\vxi)\ip{\nabla
        \plainphi(\vxi)}{\vn(\vxi)}\odif{\sigma(\vxi)}
    \end{equation}
    其中\(\odif{\sigma(\vxi)}\)表示关于“表面测度”的积分，即\(\partial \cK\)（即\(\cK\)的边界）上的诱导测度，相应地\(\vxi\)在该表面上取值，\(\vn(\vxi)\)是在表面点\(\vxi\)处指向\(\cK\)外部的单位法向量。现在，取\(\plainphi(\vxi) \doteq p_{t}(\vxi)\)和\(\psi(\vxi) \doteq c + \log p_{t}(\vxi)\)，在以\(\vzero\)为中心、半径为\(r > 0\)的球\(B_{r}(\vzero)\)上（因此\(\partial B_{r}(\vzero)\)是以\(\vzero\)为中心、半径为\(r\)的球面，且\(\vn(\vxi) = \vxi/\norm{\vxi}_{2} = \vxi/r\)）：
    \begin{align}
        &\int_{B_{r}(\vzero)}\bc{\Delta p_{t}(\vxi)[c + \log
            p_{t}(\vxi)] + \ip{\nabla \log p_{t}(\vxi)}{\nabla
        p_{t}(\vxi)}}\odif{\vxi} \\
        &= \int_{\partial B_{r}(\vzero)}[c + \log
        p_{t}(\vxi)]\ip*{\nabla
        p_{t}(\vxi)}{\frac{\vxi}{r}}\odif{\sigma(\vxi)} \\
        &= \frac{1}{r}\int_{\partial B_{r}(\vzero)}[c + \log
        p_{t}(\vxi)]\ip*{\nabla p_{t}(\vxi)}{\vxi}\odif{\sigma(\vxi)}.
    \end{align}
    令\(r \to \infty\)，有
    \begin{align}
        &\int_{\R^{D}}\bc{\Delta p_{t}(\vxi)[c + \log p_{t}(\vxi)] +
        \ip{\nabla \log p_{t}(\vxi)}{\nabla p_{t}(\vxi)}}\odif{\vxi}
        \label{eq:app_diffusion_r_infinity_limit} \\
        &= \lim_{r \to \infty}\int_{B_{r}(\vzero)}\bc{\Delta
            p_{t}(\vxi)[c + \log p_{t}(\vxi)] + \ip{\nabla \log
        p_{t}(\vxi)}{\nabla p_{t}(\vxi)}}\odif{\vxi} \\
        &= \lim_{r \to \infty}\frac{1}{r}\int_{\partial
        B_{r}(\vzero)}[c + \log p_{t}(\vxi)]\ip*{\nabla
        p_{t}(\vxi)}{\vxi}\odif{\sigma(\vxi)},
    \end{align}
    其中第一步由被积函数的控制收敛定理得出。剩下的是计算最后一个极限。为此，我们对每一项进行渐近展开。主要方法如下：对于\(\vxi \in \partial B_{r}(\vzero)\)，我们有\(\norm{\vxi}_{2} = r\)，因此
    \begin{align}
        p_{t}(\vxi)
        &= \Ex[\phi_{t}(\vxi - \vx)] \\
        &= \Ex\rs{\underbrace{\frac{1}{(2\pi)^{D/2}t^{D}}}_{\doteq
        C_{t}}e^{-\|\vxi - \vx\|_{2}^{2}/(2t^{2})}} \\
        &= C_{t}\Ex\rs{e^{-(\norm{\vxi}_{2}^{2} - 2\ip{\vxi}{\vx} +
        \norm{\vx}_{2}^{2})/(2t^{2})}} \\
        &= C_{t}\Ex\rs{e^{-(r^{2} - 2\ip{\vxi}{\vx} +
        \norm{\vx}_{2}^{2})/(2t^{2})}} \\
        &= C_{t}e^{-r^{2}/(2t^{2})}\Ex[e^{(2\ip{\vxi}{\vx} -
        \norm{\vx}_{2}^{2})/(2t^{2})}].
    \end{align}
    注意，因为\(\norm{\vxi}_{2} = r\)，由柯西-施瓦茨（Cauchy-Schwarz）不等式我们有
    \begin{equation}
        -2r\norm{\vx}_{2} - \norm{\vx}_{2}^{2} \leq 2\ip{\vxi}{\vx} -
        \norm{\vx}_{2}^{2} \leq 2r\norm{\vx}_{2} - \norm{\vx}_{2}^{2}.
    \end{equation}
    回想一下，由\Cref{assumption:entropy_x_compact_support}，\(\vx\)的支撑集是半径为\(R\)的紧集\(\cS\)。因此
    \begin{equation}
        -2R(r + R) \leq 2\ip{\vxi}{\vx} - \norm{\vx}_{2}^{2} \leq 2Rr.
    \end{equation}
    换句话说，有
    \begin{equation}
        C_{t}e^{-[r^{2} + 2R(r + R)]/(2t^{2})} \leq p_{t}(\vxi) \leq
        C_{t}e^{[-r^{2} + 2Rr]/(2t^{2})}.
    \end{equation}
    现在为了计算梯度，我们可以使用\Cref{prop:p_t_derivatives}和期望的线性性质来计算
    \begin{align}
        \ip{\nabla p_{t}(\vxi)}{\vxi}
        &= \ip*{-\frac{1}{t^{2}}\Ex\rs{\bp{\vxi - \vx}\phi_{t}(\vxi -
        \vx)}}{\vxi} \\
        &= -\frac{1}{t^{2}}\Ex\rs{\ip{\vxi - \vx}{\vxi}\phi_{t}(\vxi - \vx)} \\
        &= -\frac{1}{t^{2}}\Ex\rs{\bp{\norm{\vxi}_{2}^{2} -
        \ip{\vxi}{\vx}}\phi_{t}(\vxi - \vx)} \\
        &= -\frac{1}{t^{2}}\Ex\rs{\bp{r^{2} -
        \ip{\vxi}{\vx}}\phi_{t}(\vxi - \vx)} \\
        &= \frac{1}{t^{2}}\Ex\rs{\bp{\ip{\vxi}{\vx} -
        r^{2}}\phi_{t}(\vxi - \vx)}.
    \end{align}
    再次使用柯西-施瓦茨不等式和表示\(p_{t}(\vxi) \doteq \Ex[\phi_{t}(\vxi - \vx)]\)，有
    \begin{align}
        &\frac{1}{t^{2}}\Ex\rs{\bp{-Rr - r^{2}}\phi_{t}(\vxi - \vx)}
        \leq \ip{\nabla p_{t}(\vxi)}{\vxi} \leq
        \frac{1}{t^{2}}\Ex\rs{\bp{Rr - r^{2}}\phi_{t}(\vxi - \vx)} \\
        &\frac{1}{t^{2}}\bp{-Rr - r^{2}}\Ex\rs{\phi_{t}(\vxi - \vx)}
        \leq \ip{\nabla p_{t}(\vxi)}{\vxi} \leq \frac{1}{t^{2}}\bp{Rr
        - r^{2}}\Ex\rs{\phi_{t}(\vxi - \vx)} \\
        &-\frac{r(R + r)}{t^{2}}p_{t}(\vxi) \leq \ip{\nabla
        p_{t}(\vxi)}{\vxi} \leq -\frac{r(r - R)}{t^{2}}p_{t}(\vxi).
    \end{align}
    对于\(r > R > 0\)（这是合适的，因为我们要在固定\(R\)的同时取极限\(r \to \infty\)），两边都是负的。这是合理的：大部分概率质量包含在半径为\(R\)的球内，因此得分（score）指向内部，与向外指向的向量\(\vxi\)的内积为负。因此，使用\(p_{t}(\vxi)\)的适当界限，
    \begin{equation}
        -\frac{r(R + r)}{t^{2}}\cdot C_{t}e^{[-r^{2} + 2Rr]/(2t^{2})}
        \leq \ip{\nabla p_{t}(\vxi)}{\vxi} \leq -\frac{r(r -
        R)}{t^{2}}\cdot C_{t}e^{-[r^{2} + 2R(r + R)]/(2t^{2})}.
    \end{equation}
    然后，注意到\(C_{t} = \mathrm{poly}(t^{-1})\)，我们可以计算
    \begin{equation}
        [c + \log p_{t}(\vxi)]\ip{\nabla p_{t}(\vxi)}{\vxi} =
        \mathrm{poly}(r, R, t^{-1}, c)e^{-\Theta_{r}(r^{2})}
    \end{equation}
    因此可以看出，设\(\partial B_{r}(\vzero)\)的表面积为\(\omega_{D} r^{D - 1}\)，其中\(\omega_{D}\)是\(D\)的函数，有
    \begin{equation}
        \frac{1}{r}\int_{\partial B_{r}(\vzero)}[c + \log
        p_{t}(\vxi)]\ip{\nabla p_{t}(\vxi)}{\vxi}\odif{\vxi} =
        \mathrm{poly}(r, R, t^{-1}, c)e^{-\Theta_{r}(r^{2})}
    \end{equation}
    因此，指数衰减的尾部意味着
    \begin{equation}
        \lim_{r \to \infty}\frac{1}{r}\int_{\partial B_{r}(\vzero)}[c
        + \log p_{t}(\vxi)]\ip{\nabla p_{t}(\vxi)}{\vxi}\odif{\vxi} = 0.
    \end{equation}
    最后，代入\eqref{eq:app_diffusion_r_infinity_limit}，我们有
    \begin{align}
        &\int_{\R^{D}}\bc{\Delta p_{t}(\vxi)[c + \log p_{t}(\vxi)] +
        \ip{\nabla \log p_{t}(\vxi)}{\nabla p_{t}(\vxi)}}\odif{\vxi} = 0 \\
        \implies
        &\int_{\R^{D}}\Delta p_{t}(\vxi)[c + \log
        p_{t}(\vxi)]\odif{\vxi} = -\int_{\R^{D}}\ip{\nabla \log
        p_{t}(\vxi)}{\nabla p_{t}(\vxi)}\odif{\vxi}
    \end{align}
    如所断言。
\end{proof}

\subsubsection{去噪器\(\bar{\vx}\)的局部可逆性}

在这里，我们提供在\Cref{thm:conditioning_reduces_entropy}证明中使用的一些结果，这些结果是\cite{Gribonval2011-pf}中相应结果的适当推广。

\begin{lemma}[\cite{Gribonval2011-pf}引理A.1的推广]\label{lem:gribonval_A1}
    设\(\vx\)为满足\Cref{assumption:entropy_x_compact_support,assumption:entropy_x_density}的任意随机变量，并设\((\vx_{t})_{t \in [0, T]}\)为随机过程\eqref{eq:app_additive_gaussian_noise_model}。设\(s, t \in [0, T]\)满足\(0 \leq s < t \leq T\)，并设\(\bar{\vx}(\vxi) \doteq \Ex[\vx_{s} \mid \vx_{t} = \vxi]\)。雅可比矩阵\(\bar{\vx}^{\prime}(\vxi)\)是对称正定的。
\end{lemma}
\begin{proof}
    我们有
    \begin{equation}
        \bar{\vx}^{\prime}(\vxi) = \vI + \bp{1 -
        \frac{s}{t}}t^{2}\nabla^{2}\log p_{t}(\vxi).
    \end{equation}
    这里我们展开
    \begin{equation}
        \nabla^{2}\log p_{t}(\vxi) =
        \frac{p_{t}(\vxi)\nabla^{2}p_{t}(\vxi) - (\nabla
        p_{t}(\vxi))(\nabla p_{t}(\vxi))^{\top}}{p_{t}(\vxi)^{2}}.
    \end{equation}
    因此我们需要确保
    \begin{align}
        \bar{\vx}^{\prime}(\vxi)
        &= \vI + \bp{1 -
        \frac{s}{t}}t^{2}\frac{p_{t}(\vxi)\nabla^{2}p_{t}(\vxi) -
        (\nabla p_{t}(\vxi))(\nabla p_{t}(\vxi))^{\top}}{p_{t}(\vxi)^{2}} \\
        &= \frac{p_{t}(\vxi)^{2}\vI + \bp{1 -
            \frac{s}{t}}t^{2}\bs{p_{t}(\vxi)\nabla^{2}p_{t}(\vxi) -
        (\nabla p_{t}(\vxi))(\nabla p_{t}(\vxi))^{\top}}}{p_{t}(\vxi)^{2}}
    \end{align}
    是对称半正定的。实际上它显然是对称的（由克莱罗定理（Clairaut's theorem）可知）。为了证明其半正定性，我们代入由\eqref{eq:p_t_representation}给出的\(p_{t}\)的期望表示（以及由\Cref{prop:p_t_derivatives}给出的\(\nabla p_{t}\)、\(\Delta p_{t}\)），得到（其中\(\vx\)如前所定义，\(\vy\)与\(\vx\)独立同分布），
    \begin{align}
        &\vv^{\top}[\bar{\vx}^{\prime}(\vxi)]\vv \\
        &= p_{t}(\vxi)^{-2}\vv^{\top}\Bigg\{p_{t}(\vxi)^{2}\vI +
            \bp{1 - \frac{s}{t}}t^{2}\Ex[\phi_{t}(\vxi -
            \vx)]\Ex\rs{\phi_{t}(\vxi - \vx)\cdot\frac{(\vxi -
            \vx)(\vxi - \vx)^{\top} - t^{2}\vI}{t^{4}}} \\
            & \qquad \qquad \qquad -\bp{1 -
            \frac{s}{t}}t^{2}\Ex\rs{-\phi_{t}(\vxi - \vx)\cdot\frac{\vxi
            - \vx}{t^{2}}}\Ex\rs{-\phi_{t}(\vxi - \vx)\cdot\frac{\vxi -
        \vx}{t^{2}}}^{\top}\Bigg\}\vv \\
        &= p_{t}(\vxi)^{-2}\vv^{\top}\Bigg\{\Ex[\phi_{t}(\vxi -
            \vx)\phi_{t}(\vxi - \vy)\vI] \\
            & \qquad \qquad \qquad + \bp{1 -
            \frac{s}{t}}t^{2}\Ex\rs{\phi_{t}(\vxi - \vx)\phi_{t}(\vxi
                - \vy)\cdot\frac{(\vxi - \vy)(\vxi - \vy)^{\top} -
            t^{2}\vI}{t^{4}}} \\
            & \qquad \qquad \qquad -\bp{1 -
            \frac{s}{t}}t^{2}\Ex\rs{\phi_{t}(\vxi - \vx)\phi_{t}(\vxi -
        \vy)\cdot\frac{(\vxi - \vx)(\vxi - \vy)^{\top}}{t^{4}}}\Bigg\} \\
        &= \frac{1 -
        s/t}{p_{t}(\vxi)^{2}}\vv^{\top}\Ex\mathopen{}\Bigg[\phi_{t}(\vxi
            - \vx)\phi_{t}(\vxi - \vy)\bc{\frac{1}{1 - s/t}\vI +
                \frac{(\vxi - \vy)(\vxi - \vy)^{\top}}{t^{2}} - \vI -
        \frac{(\vxi - \vx)(\vxi - \vy)^{\top}}{t^{2}}}\Bigg]\vv \\
        &= \frac{t - s}{tp_{t}(\vxi)^{2}}\vv^{\top}\Ex\rs{\frac{s}{t
            - s}\vI +\frac{(\vxi - \vx)(\vxi - \vx)^{\top}}{2t^{2}} +
            \frac{(\vxi - \vy)(\vxi - \vy)^{\top}}{2t^{2}} - \frac{(\vxi
        - \vx)(\vxi - \vy)^{\top}}{t^{2}}}\vv \\
        &= \frac{t - s}{tp_{t}(\vxi)^{2}}\vv^{\top}\Ex\rs{\frac{s}{t
            - s}\vI +\frac{1}{2t^{2}}\bp{(\vxi - \vx)(\vxi - \vx)^{\top}
                + (\vxi - \vy)(\vxi - \vy)^{\top} - 2(\vxi - \vx)(\vxi -
        \vy)^{\top}}}\vv \\
        &= \frac{t - s}{tp_{t}(\vxi)^{2}}\Ex\rs{\frac{s}{t -
            s}\norm{\vv}_{2}^{2} +\frac{1}{2t^{2}}\bp{[(\vxi -
                \vx)^{\top}\vv]^{2} + [(\vxi - \vy)^{\top}\vv]^{2} - 2[(\vxi
        - \vx)^{\top}\vv][(\vxi - \vy)^{\top}\vv]}} \\
        &= \frac{t - s}{tp_{t}(\vxi)^{2}}\Ex\rs{\frac{s}{t -
            s}\norm{\vv}_{2}^{2} +\frac{1}{2t^{2}}\bp{[(\vxi -
                \vx)^{\top}\vv]^{2} + [(\vxi - \vy)^{\top}\vv]^{2} - 2[(\vxi
        - \vx)^{\top}\vv][(\vxi - \vy)^{\top}\vv]}} \\
        &= \frac{t - s}{tp_{t}(\vxi)^{2}}\Ex\rs{\frac{s}{t -
        s}\norm{\vv}_{2}^{2} +\frac{1}{2t^{2}}\bp{[(\vxi -
        \vx)^{\top}\vv] - [(\vxi - \vy)^{\top}\vv]}^{2}} \\
        &= \frac{t - s}{tp_{t}(\vxi)^{2}}\Ex\rs{\frac{s}{t -
        s}\norm{\vv}_{2}^{2} +\frac{1}{2t^{2}}[(\vy - \vx)^{\top}\vv]^{2}} \\
        &= \frac{s}{tp_{t}(\vxi)^{2}}\norm{\vv}_{2}^{2} + \frac{t -
        s}{2t^{3}p_{t}(\vxi)}\Ex[\{(\vy - \vx)^{\top}\vv\}^{2}]
    \end{align}
    由于\(\vx\)和\(\vy\)是独立同分布的，整个积分（即原始二次型）为\(0\)当且仅当\(s = 0\)且\(\vx\)的支撑集完全包含在一个与\(\vv\)正交的仿射子空间中。但这被假设（即\(\vx\)在\(\R^{D}\)上具有密度）排除了，因此雅可比矩阵\(\bar{\vx}^{\prime}(\vxi)\)是对称正定的。
\end{proof}

\begin{lemma}[\cite{Gribonval2011-pf}推论A.2第1部分的推广]\label{lem:gribonval_A2}
    设\(f \colon \R^{D} \to \R^{D}\)为任意可微函数，其雅可比矩阵\(f^{\prime}(\vx)\)是对称正定的。那么\(f\)是单射的，因此作为函数\(\R^{D} \to \Range(f)\)是可逆的，其中\(\Range(f)\)是\(f\)的值域。
\end{lemma}
\begin{proof}
    假设\(f\)不是单射的，即存在\(\vx, \vx^{\prime}\)使得\(f(\vx) = f(\vx^{\prime})\)但\(\vx \neq \vx^{\prime}\)。定义\(\vv \doteq (\vx^{\prime} - \vx)/\norm{\vx^{\prime} - \vx}_{2}\)。定义函数\(g \colon \R \to \R\)为\(g(t) \doteq \vv^{\top}f(\vx + t\vv)\)。那么\(g(0) = \vv^{\top}f(\vx) = \vv^{\top}f(\vx^{\prime}) = g(\norm{\vx^{\prime} - \vx}_{2})\)。由于\(f\)是可微的，\(g\)也是可微的，因此由中值定理可知，导数\(g^{\prime}\)必须在某个\(t^{\ast} \in (0, \norm{\vx^{\prime} - \vx}_{2})\)处为零。然而，
    \begin{equation}
        g^{\prime}(t^{\ast}) \doteq \vv^{\top}\bs{f^{\prime}(\vx +
        t^{\ast}\vv)}\vv > 0
    \end{equation}
    因为雅可比矩阵是正定的。因此我们得出了矛盾，如所断言。
\end{proof}

结合上述两个结果，我们得到以下关键结果。

\begin{corollary}[\cite{Gribonval2011-pf}推论A.2第2部分的推广]\label{cor:gribonval_A2}
    设\(\vx\)为满足\Cref{assumption:entropy_x_compact_support,assumption:entropy_x_density}的任意随机变量，并设\((\vx_{t})_{t \in [0, T]}\)为随机过程\eqref{eq:app_additive_gaussian_noise_model}。设\(s, t \in [0, T]\)满足\(0 \leq s < t \leq T\)，并设\(\bar{\vx}(\vxi) \doteq \Ex[\vx_{s} \mid \vx_{t} = \vxi]\)。那么\(\bar{\vx}\)是单射的，因此在其值域上是可逆的。
\end{corollary}
\begin{proof}
    剩下唯一要证明的是\(\bar{\vx}\)是可微的，但这可以从Tweedie公式（\Cref{thm:tweedie}）直接得出，该公式表明\(\bar{\vx}\)是可微的当且仅当\(\nabla \log p_{t}\)是可微的，而这由\Cref{eq:p_t_representation}提供。
\end{proof}

\subsubsection{控制拉普拉斯算子\(\Delta \log p_{t}\)}
最后，我们建立一个技术性估计，这是证明\Cref{thm:conditioning_reduces_entropy}所需的，并且实际上激发了对可行\(t\)的假设。

\begin{lemma}\label{lem:app_diffusion_laplacian_control}
    设 \(\vx\) 为满足 \Cref{assumption:entropy_x_compact_support,assumption:entropy_x_density} 的任意随机变量，并设 \((\vx_{t})_{t \in [0, T]}\) 为随机过程 \eqref{eq:app_additive_gaussian_noise_model}。设 \(p_{t}\) 为 \(\vx_{t}\) 的密度。那么，对于 \(t > 0\)，有
    \begin{equation}
        \sup_{\vxi \in \R^{D}}\abs{\Delta \log p_{t}(\vxi)} \leq
        \max\rp{\frac{D}{t^{2}}, \abs*{\frac{R^{2}}{t^{4}} - \frac{D}{t^{2}}}}.
    \end{equation}
\end{lemma}
\begin{proof}
    根据链式法则，通过简单的计算可得
    \begin{equation}
        \Delta \log p_{t}(\vxi) = \frac{\Delta
        p_{t}(\vxi)}{p_{t}(\vxi)} - \frac{\norm{\nabla
        p_{t}(\vxi)}_{2}^{2}}{p_{t}(\vxi)^{2}}.
    \end{equation}
    使用 \Cref{prop:p_t_derivatives} 展开 \(\Delta p_{t}(\vxi)\) 中的项，我们得到
    \begin{align}
        \frac{\Delta p_{t}(\vxi)}{p_{t}(\vxi)}
        &= \frac{\Ex\rs{\frac{\norm{\vxi - \vx}_{2}^{2} -
        Dt^{2}}{t^{4}} \cdot \phi_{t}(\vxi - \vx)}}{\Ex[\phi_{t}(\vxi
        - \vx)]} \\
        &= \frac{\int_{\R^{D}}\bc{\frac{\norm{\vxi - \vu}_{2}^{2} -
            Dt^{2}}{t^{4}}}\phi_{t}(\vxi -
        \vu)p(\vu)\odif{\vu}}{\int_{\R^{D}}\phi_{t}(\vxi -
        \vu)p(\vu)\odif{\vu}}.
    \end{align}
    这看起来像是一个贝叶斯边缘化过程，因此让我们定义适当的归一化密度
    \begin{equation}
        q_{\vxi}(\vu) = \frac{\phi_{t}(\vxi -
        \vu)p(\vu)}{\int_{\R^{D}}\phi_{t}(\vxi -
        \vv)p(\vv)\odif{\vv}} = \frac{\phi_{t}(\vxi -
        \vu)p(\vu)}{\Ex[\phi_{t}(\vxi - \vx)]} = \frac{\phi_{t}(\vxi
        - \vu)p(\vu)}{p_{t}(\vxi)}
    \end{equation}
    然后，定义 \(\vy_{\vxi} \sim q_{\vxi}\)，我们可以写出
    \begin{equation}
        \frac{\Delta p_{t}(\vxi)}{p_{t}(\vxi)} =
        \int_{\R^{D}}\bc{\frac{\norm{\vxi - \vu}_{2}^{2} -
        Dt^{2}}{t^{4}}}q_{\vxi}(\vu)\odif{\vu} =
        \frac{1}{t^{4}}\Ex[\norm{\vxi - \vy_{\vxi}}_{2}^{2}] - \frac{D}{t^{2}}.
    \end{equation}
    类似地，写出第二项（非平方项），我们得到
    \begin{equation}
        \frac{\nabla p_{t}(\vxi)}{p_{t}(\vxi)} = -\frac{\vxi -
        \Ex[\vy_{\vxi}]}{t^{2}}.
    \end{equation}
    令 \(\vz_{\vxi} \doteq \vy_{\vxi} - \vxi\)，有
    \begin{equation}
        \frac{\Delta p_{t}(\vxi)}{p_{t}(\vxi)} =
        \frac{\Ex[\norm{\vz_{\vxi}}_{2}^{2}]}{t^{4}} -
        \frac{D}{t^{2}}, \qquad \frac{\nabla
        p_{t}(\vxi)}{p_{t}(\vxi)} = \frac{\Ex[\vz_{\vxi}]}{t^{2}}.
    \end{equation}
    因此，将 \(\Delta \log p_{t}\) 完整写出，我们有
    \begin{align}
        \Delta \log p_{t}(\vxi)
        &= \frac{\Ex[\norm{\vz_{\vxi}}_{2}^{2}]}{t^{4}} -
        \frac{D}{t^{2}} - \frac{\norm{\Ex[\vz_{\vxi}]}_{2}^{2}}{t^{4}} \\
        &= \frac{\Ex[\norm{\vz_{\vxi}}_{2}^{2}] -
        \norm{\Ex[\vz_{\vxi}]}_{2}^{2}}{t^{4}} - \frac{D}{t^{2}} \\
        &= \frac{\tr(\Cov(\vz_{\vxi}))}{t^{4}} - \frac{D}{t^{2}} \\
        &= \frac{\tr(\Cov(\vy_{\vxi}))}{t^{4}} - \frac{D}{t^{2}}.
    \end{align}
    该迹的一个平凡下界是 \(0\)，因为协方差矩阵是半正定的。为了寻找上界，注意到 \(\vy_{\vxi}\) 仅在 \(\vx\) 的支撑集上取值（因为 \(p\) 是 \(\vy_{\vxi}\) 的密度 \(q_{\vxi}\) 的一个因子），根据 \Cref{assumption:entropy_x_compact_support}，该支撑集是一个紧集 \(\cS\)，其半径为 \(R \doteq \sup_{\vxi \in \cS}\norm{\vxi}_{2}\)。因此
    \begin{equation}
        \tr(\Cov(\vy_{\vxi})) = \Ex[\norm{\vy_{\vxi}}_{2}^{2}] -
        \norm{\Ex[\vy_{\vxi}]}_{2}^{2} \leq
        \Ex[\norm{\vy_{\vxi}}_{2}^{2}] \leq R^{2}.
    \end{equation}
    因此
    \begin{equation}
        -\frac{D}{t^{2}} \leq \Delta \log p_{t}(\vxi) \leq
        \frac{R^{2}}{t^{4}} - \frac{D}{t^{2}},
    \end{equation}
    这证明了该结论。
\end{proof}

\subsubsection{导数计算}

在这里，我们计算一些有用的导数，这些导数将在整个附录中被重复使用。

\begin{proposition}\label{prop:p_t_derivatives}
    设 \(\vx\) 为满足 \Cref{assumption:entropy_x_compact_support,assumption:entropy_x_density} 的任意随机变量，并设 \((\vx_{t})_{t \in [0, T]}\) 为随机过程 \eqref{eq:app_additive_gaussian_noise_model}。对于 \(t \geq 0\)，设 \(p_{t}\) 为 \(\vx_{t}\) 的密度。那么
    \begin{align}
        \pdv{p_{t}}{t}(\vxi)
        &= \Ex\rs{\phi_{t}(\vxi - \vx)\cdot\frac{\norm{\vxi -
        \vx}_{2}^{2} - Dt^{2}}{t^{3}}} \\
        \nabla p_{t}(\vxi)
        &= -\Ex\rs{\phi_{t}(\vxi - \vx)\cdot\frac{\vxi - \vx}{t^{2}}} \\
        \nabla^{2}p_{t}(\vxi)
        &= \Ex\rs{\phi_{t}(\vxi - \vx)\cdot\frac{(\vxi - \vx)(\vxi -
        \vx)^{\top} - t^{2}\vI}{t^{4}}} \\
        \Delta p_{t}(\vxi)
        &= \Ex\rs{\phi_{t}(\vxi - \vx)\cdot\frac{\norm{\vxi -
        \vx}_{2}^{2} - Dt^{2}}{t^{4}}}.
    \end{align}
\end{proposition}
\begin{proof}
    我们使用 \(p_{t}\) 的卷积表示，即 \eqref{eq:p_t_representation}。首先对时间求导，计算表明 \Cref{prop:dutis} 是适用的，\footnote{我们使用 \(f_{t}(\vxi) = p(\vxi) \phi_{t}(\vxi - \vx)\)，注意到它关于 \(\vxi\) 是二次连续可微的，并且关于 \(t\) 是（超过）二次连续可微的。然后为了检查 \(f_{t}\) 的局部可积性，我们计算 \(\pdv{f_{t}}{t}(\vxi) = f_{t}(\vxi)\cdot \frac{1}{t^{3}}(\norm{\vxi - \vx}_{2}^{2} - Dt^{2})\)，很容易验证其在 \(\vxi\) 和 \(t \in [t_{\min}, t_{\max}]\)（其中 \(t_{\min} > 0\)）上是可积的。（实际上，\(f_{t}\) 具有指数衰减的尾部，因此乘积中的二次项不会造成问题。）} 因此我们可以将导数移入积分/期望内部：
    \begin{equation}
        \pdv{p_{t}}{t}(\vxi) = \pdv*{\Ex[\phi_{t}(\vxi - \vx)]}{t} =
        \Ex\rs{\pdv*{\phi_{t}(\vxi - \vx)}{t}} = \pdv{\phi_{t}}{t} * p.
    \end{equation}
    同时，根据卷积的性质（\Cref{prop:diff_convolution}）并利用 \(p\) 具有紧支撑的事实（\Cref{assumption:entropy_x_compact_support}），
    \begin{equation}
        p_{t} = \phi_{t} * p \implies \nabla p_{t} = \nabla \phi_{t}
        * p \implies \nabla^{2}p_{t} = \nabla^{2}\phi_{t} * p
        \implies \Delta p_{t} = \Delta \phi_{t} * p.
    \end{equation}
    剩余的计算可由 \Cref{prop:normal_derivatives} 得出。
\end{proof}

\begin{proposition}\label{prop:normal_derivatives}
    对于 \(t > 0\) 和 \(\vxi \in \R^{D}\)，有
    \begin{align}
        \pdv*{\phi_{t}}{t}(\vxi)
        &= \phi_{t}(\vxi) \cdot \frac{\norm{\vxi}_{2}^{2} - Dt^{2}}{t^{3}} \\
        \nabla \phi_{t}(\vxi)
        &= -\phi_{t}(\vxi)\cdot\frac{\vxi}{t^{2}} \\
        \nabla^{2} \phi_{t}(\vxi)
        &= \phi_{t}(\vxi)\cdot\frac{\vxi\vxi^{\top} - t^{2}\vI}{t^{4}} \\
        \Delta \phi_{t}(\vxi)
        &= \phi_{t}(\vxi) \cdot \frac{\norm{\vxi}_{2}^{2} - Dt^{2}}{t^{4}}.
    \end{align}
\end{proposition}
\begin{proof}
    直接计算。
\end{proof}

\subsubsection{积分号下求导}

在本附录中，我们多次在积分号下求导，了解何时可以这样做非常重要。积分号下求导有两种情况：
\begin{enumerate}
    \item 对积分 \(\int f_{t}(\vxi)\odif{\vxi}\) 关于辅助参数 \(t\) 求导。
    \item 对卷积 \((f * g)(\vxi) = \int f(\vu)g(\vxi - \vu)\odif{\vu}\) 关于变量 \(\vxi\) 求导。
\end{enumerate}

对于第一种情况，我们给出一个具体的结果，此处不加证明地陈述，但可归因于 \href{https://planetmath.org/differentiationundertheintegralsign}{链接的来源}，该来源将以下结果作为关于微分算子与缓增分布相互作用的更一般定理的特例推导出来，这远远超出了本书的范围。完整的正式参考文献可见 \cite{jones1982theory}。
\begin{proposition}[\cite{jones1982theory}, 第 11.12 节]\label{prop:dutis}
    设 \(f \colon (0, T) \times \R^{D} \to \R\) 满足：
    \begin{itemize}
        \item \(f\) 是关于 \((t, \vxi)\) 的联合可测函数；
        \item 对于勒贝格几乎所有的 \(\vxi \in \R^{D}\)，函数 \(t \mapsto f_{t}(\vxi)\) 是绝对连续的；
        \item \(\pdv{f_{t}}{t}\) 是局部可积的，即对于每个 \([t_{\min}, t_{\max}] \subseteq (0, T)\)，有
            \begin{equation}
                \int_{t_{\min}}^{t_{\max}}\int_{\R^{D}}\abs*{\pdv{f_{t}}{t}(\vxi)}\odif{\vxi}
                < \infty.
            \end{equation}
    \end{itemize}
    那么 \(t \mapsto \int_{\R^{D}}f_{t}(\vxi)\odif{\vxi}\) 是 \((0, T)\) 上的绝对连续函数，且其导数为
    \begin{equation}
        \odv*{\int_{\R^{D}}f_{t}(\vxi)\odif{\vxi}}{t} =
        \int_{\R^{D}}\pdv*{f_{t}}{t}(\vxi)\odif{\vxi},
    \end{equation}
    对几乎所有的 \(t \in (0, T)\) 都有定义。
\end{proposition}

对于第二种情况，我们给出另一个具体的结果，此处不加证明地陈述，但在 \cite{brezis2011functional} 中有完整的形式化证明。
\begin{proposition}[\cite{brezis2011functional}, 命题 4.20]\label{prop:diff_convolution}
    设 \(f\) 是具有紧支撑的 \(k\) 次连续可微函数，且设 \(g\) 是局部可积的。那么由下式定义的卷积 \(f * g\)
    \begin{equation}
        (f * g)(\vxi) \doteq \int_{\R^{D}}f(\vu)g(\vxi - \vu)\odif{\vu}
    \end{equation}
    是 \(k\) 次连续可微的，且其 \(k\) 阶导数为
    \begin{equation}
        \nabla^{k}(f * g) =(\nabla^{k}f) * g.
    \end{equation}
\end{proposition}
尽管书中未提及，但一个简单的分部积分论证表明，如果 \(g\) 也是 \(k\) 次可微的，那么我们可以“权衡”这种正则性：
\begin{equation}
    \nabla^{k}(f * g) = f * (\nabla^{k} g).
\end{equation}

\section{有损编码与球堆积}\label{app:rate-distortion-covering}

在本节中，我们证明 \Cref{thm:covering-number-rate-distortion}。遵循贯穿本附录的约定，我们将随机变量 $\vx$ 的紧支撑集记为 $\cS = \Supp(\vx)$。

正如前面所预示的，我们将对支撑集 $\cS$ 做出正则性假设以证明 \Cref{thm:covering-number-rate-distortion}。在最小假设下进行推导的一种可能性是，在我们的设定中实例化 \cite{Riegler2018-jh,Riegler2023-rr} 的结果，因为这些结果适用于正则性非常低的集合 $\cS$（例如，具有分形结构的类康托尔集）。然而，我们发现在我们的设定中精确计算这些结果中的常数（这是断言如 \Cref{thm:covering-number-rate-distortion} 这样的结论所必需的努力）有些繁琐。
因此，我们的方法是在集合 $\cS$ 上增加一个几何正则性假设，这牺牲了一些一般性，但允许我们展开更透明的论证。为了避免牺牲过多的一般性，我们必须确保不排除集合 $\cS$ 的低维性。
因此，我们考虑贯穿本书的运行示例，即低秩高斯分布的混合。在这种几何设定下，我们将通过假设 $\cS$ 是超球面的并集来强制实现这一点，这在极大概率上等价于高维情况下的高斯假设。

\begin{assumption}\label{assumption:union-of-spheres}
    随机变量 $\vx$ 的支撑集 $\cS \subset \R^D$ 是 $K$ 个球面的有限并集，每个球面的维度为 $d_k$，$k \in [K]$。
    $\vx$ 从第 $k$ 个球面中抽取的概率由 $\pi_k \in [0, 1]$ 给出，并且在给定从第 $k$ 个球面抽取的条件下，$\vx$ 在该球面上均匀分布。
    这些支撑集满足每个球面与所有其他球面相互正交。
\end{assumption}

我们在简化的 \Cref{assumption:union-of-spheres} 下进行推导，以简化过多的技术细节，并与贯穿本专著的重要运行示例联系起来。我们相信，通过额外的技术努力，我们的结果可以推广到属于\textit{具有正向可达性的集合}（sets with positive reach）类的支撑集 $\cS$，但我们将此留待未来研究。

\subsection{率失真与覆盖之间关系的证明}

我们简要概述证明思路，然后着手建立三个基本引理，最后给出证明。该证明将依赖于下面概述中引入的概念。

获得率失真函数 \eqref{eqn:rate-distortion-general} 的上界是直接的：根据率的特征（即，率失真函数是 $\vx$ 的期望平方 $\ell^2$ 失真为 $\epsilon$ 的编码的最小率），对 $\cR_{\epsilon}(\x)$ 求上界只需要展示一种达到该目标失真的 $\x$ 的编码，而 $\Supp(\x)$ 的任何 $\epsilon$-覆盖都能实现这一点，其率等于该覆盖基数的以 $2$ 为底的对数。
下界则更为微妙。我们利用在 \Cref{rem:slb} 中讨论的香农下界：计算 \cite[\S III, (22)]{Linder1994-ej} 中的常数，可以得到 \Cref{eq:slb} 中引用结果的更精确版本（当然，单位是比特）：对于任何具有紧支撑和密度的随机变量 $\vx$，有
\begin{equation}\label{eq:slb-appendix-1}
    \cR_{\epsilon}(\x)
    \geq
    h(\x)
    - \log \volume(B_{\epsilon})
    +
    \log
    \left(
        \frac{
            2
        }
        {
            D \Gamma(D/2)
        }
        \left(
            \frac{
                D
            }
            {
                2e
            }
        \right)^{D/2}
    \right),
\end{equation}
在此表达式中，熵（等）的单位为奈特（nats）。
该常数可以使用斯特林近似（Stirling's approximation）轻松估计。
斯特林近似的一种通常很有用的定量形式给出，对于任何 $x > 0$ \cite{Jameson2015-hy}
\begin{equation}
    \Gamma(x) \leq \sqrt{2\pi} x^{x - 1/2} e^{-x} e^{1/(12x)}.
\end{equation}
我们将把这个界应用于 \Cref{eq:slb-appendix-1} 中的 $\Gamma(D/2)$。
我们得到
\begin{align}
    \log
    \left(
        \frac{
            2
        }
        {
            D \Gamma(D/2)
        }
        \left(
            \frac{
                D
            }
            {
                2e
            }
        \right)^{D/2}
    \right)
    &\geq
    -\frac{1}{6D}
    +
    \log\left(
        \frac{2}{D}
        \left(
            \frac{
                D
            }
            {
                2e
            }
        \right)^{D/2}
        \cdot
        \sqrt{\frac{D}{4\pi}}
        \left(
            \frac{D}{2e}
        \right)^{-D/2}
    \right)
    \\
    &=
    -\frac{1}{6D}
    - \frac{1}{2}\log D\pi,\label{eq:slb-constant-est}
\end{align}
我们可以将其作为 \Cref{eq:slb} 中常数 $C_D$ 的显式值。
总结完全量化的香农下界（单位为比特）：
\begin{equation}\label{eq:slb-appendix}
    \cR_{\epsilon}(\x)
    \geq
    h(\x)
    - \log_2 \volume(B_{\epsilon})
    - O(\log D).
\end{equation}

现在，对于我们当前目的而言，一个重要的限制是香农下界要求随机变量 $\vx$ 具有密度，这排除了许多我们感兴趣的低维分布。
但让我们暂时考虑 $\vx$ 确实具有密度的情况。
$\vx$ 在其支撑集上均匀分布的假设在这种设定下很容易形式化：对于任何博雷尔集 $A \subset \cS$，我们有
\begin{equation}
    \Pr[\vx \in A] = \int_{A} \frac{1}{\volume(\cS)} \odif \vx.
\end{equation}
那么熵 $h(\vx)$ 就是
\begin{equation}
    h(\vx) = \log_2 \volume(\cS).
\end{equation}
然后，证明以一个引理结束，该引理将比率 $\volume(\cS) / \volume(B_{\epsilon})$ 与 $\cS$ 被 $\epsilon$ 球覆盖的 $\epsilon$-覆盖数联系起来。

为了将上述方案扩展到满足 \Cref{assumption:union-of-spheres} 的退化分布，
我们在 \Cref{thm:covering-number-rate-distortion} 中下界的证明将利用一个近似论证，即用具有密度的“附近”分布来近似实际的低维分布 $\vx$，这与 \Cref{thm:max_entropy} 之前的证明概述相似但不完全相同。
然后，我们将把近似序列中引入的参数与失真参数 $\epsilon$ 联系起来，以获得 \Cref{thm:covering-number-rate-distortion} 中所需的结论。

\begin{definition}\label{def:thickening-set}
    设 $\cS$ 为一个紧集。
    对于任意 $\delta > 0$，定义 $\cS$ 的 $\delta$-增厚（thickening），记为 $\cS_{\delta}$，如下：
    \begin{equation}
        \cS_{\delta} = \set*{\vxi \in \R^D \given \mathrm{dist}(\vxi, \cS) \leq
        \delta}.
    \end{equation}
\end{definition}
\Cref{def:thickening-set} 中引用的距离函数定义为
\begin{equation}\label{eq:dist-func}
    \dist(\vxi, \cS) = \inf_{\vxi' \in \cS} \norm*{\vxi - \vxi'}_2.
\end{equation}
对于紧集 $\cS$，魏尔斯特拉斯定理（Weierstrass's theorem）意味着对于任何 $\vxi \in \R^D$，总存在某个 $\vxi' \in \cS$ 能够达到距离函数中的下确界。
$\cS_{\delta}$ 的紧致性很容易由 $\cS$ 的紧致性得出，因此对于任何 $\delta > 0$，$\volume(\cS_{\delta})$ 都是有限的。于是，我们可以对增厚随机变量做出以下定义，专门针对 \Cref{assumption:union-of-spheres}。

\begin{definition}\label{def:thickening-rv-uos}
    设 $\vx$ 为一个随机变量，使得 $\Supp(\vx) = \cS$ 是 $K$ 个超球面的并集，其分布如 \Cref{assumption:union-of-spheres} 所述。
    将混合分布中每个分量的支撑集记为 $\cS_k$。
    定义增厚随机变量 $\vx_{\delta}$ 为测度的混合，其中每个分量测度在增厚集 $\cS_{k, \delta}$（\Cref{def:thickening-set}）上是均匀的，对于 $k \in [K]$，混合权重为 $\pi_k$。
\end{definition}

\begin{lemma}\label{lem:rate-distortion-lb-uos}
    假设随机变量 $\vx$ 满足 \Cref{assumption:union-of-spheres}。那么如果 $0 < \delta < \tfrac{1}{2}$，增厚随机变量 $\vx_{\delta}$（\Cref{def:thickening-rv-uos}）对于任意 $\epsilon > 0$ 满足
    \begin{equation}
        R_{\delta + \epsilon}(\vx_{\delta})
        \leq
        R_{\epsilon}(\vx).
    \end{equation}
\end{lemma}

\Cref{lem:rate-distortion-lb-uos} 的证明推迟到 \Cref{sec:app-rate-dist-deferred-proofs}。
使用 \Cref{lem:rate-distortion-lb-uos}，上述方案得以实现，因为随机变量 $\vx_{\delta}$ 具有关于勒贝格测度均匀的密度。

\begin{proof}[(\Cref{thm:covering-number-rate-distortion} 的证明)]
    上界很容易证明。如果 $S$ 是 $\vx$ 支撑集的任意 $\epsilon$-覆盖，其基数为 $\cN_{\epsilon}(\Supp(\vx))$，那么考虑将每个 $\vxi \in \Supp(\vx)$ 映射到重建值 $\hat{\vxi} = \argmin_{\vxi' \in S}\, \norm{\vxi - \vxi'}_2$ 的编码方案，平局时任意打破。那么平局发生的概率为零，并且 $S$ 在尺度 $\epsilon$ 下覆盖 $\Supp(\vx)$ 这一事实保证了失真不大于 $\epsilon$；该方案的率为 $\log_2 \cN_{\epsilon}(\Supp(\vx))$。

    对于下界，令 $0 < \delta < \tfrac{1}{2}$，并考虑增厚随机变量 $\vx_{\delta}$。根据 \Cref{lem:rate-distortion-lb-uos}，我们有
    \begin{equation}
        R_{\delta + \epsilon}(\vx_{\delta})
        \leq
        R_{\epsilon}(\vx).
    \end{equation}
    由于 $\vx_{\delta}$ 具有均匀的勒贝格密度，我们随后可以应用形式为 \eqref{eq:slb-appendix} 的香农下界，得到
    \begin{equation}\label{eq:slb-proof}
        \log_2 \volume(\Supp(\vx_{\delta}))
        - \log_2 \volume(B_{\delta + \epsilon})
        - O(\log D)
        \leq
        R_{\epsilon}(\vx).
    \end{equation}
    最后，我们需要对以下比率求下界
    \begin{equation}
        \frac{
            \volume(\Supp(\vx_{\delta}))
        }
        {
            \volume(B_{\delta + \epsilon})
        }
    \end{equation}
    用覆盖数来表示。
    由于 $\Supp(\vx_{\delta}) = \Supp(\vx) + B_{\delta}$，其中此处的 $+$ 表示闵可夫斯基和（Minkowski sum），体积界论证的标准应用（例如参见 \cite[命题 4.2.12]{Vershynin2018-br}）给出
    \begin{equation}
        \volume(\Supp(\vx_{\delta}))
        \geq
        \cN_{2\delta}(\Supp(\vx))
        \volume(B_{\delta}).
    \end{equation}
    因此
    \begin{align}
        \frac{
            \volume(\Supp(\vx_{\delta}))
        }
        {
            \volume(B_{\delta + \epsilon})
        }
        &\geq
        \cN_{2\delta}(\Supp(\vx))
        \frac{
            \volume(B_{\delta})
        }
        {
            \volume(B_{\delta + \epsilon})
        }
        \\
        &=
        \cN_{2\delta}(\Supp(\vx))
        \left(
            \frac{
                \delta
            }
            {
                \delta + \epsilon
            }
        \right)^D.
    \end{align}
    选择 $\delta = \epsilon / 2$，由香农下界 \eqref{eq:slb-proof} 和上述估计可得
    \begin{equation}
        \log_2 \cN_{\epsilon}(\Supp(\vx))
        - O(D)
        \leq
        R_{\epsilon}(\vx),
    \end{equation}
    证毕。

\end{proof}

\subsection{引理 \ref{lem:rate-distortion-lb-uos} 的证明}\label{sec:app-rate-dist-deferred-proofs}

\begin{proof}[(\Cref{lem:rate-distortion-lb-uos} 的证明)]
    只需证明，对于适当选择的 $\delta$，任何期望平方失真为 $\epsilon^2$ 的 $\vx$ 编码都能产生一个具有相同率且失真不会大太多的 $\vx_{\delta}$ 编码。
    因此，固定这样一个 $\vx$ 的编码，其达到率 $R$ 和期望平方失真 $\epsilon^2$。我们将使用该编码重建的随机变量记为 $\hat{\vx}$，并将相关的编码-解码映射记为 $\mathrm{q} : \Supp(\vx) \to \Supp(\vx)$（即 $\hat{\vx} = \mathrm{q}(\vx)$）。

    现在令 $\cS_k$ 表示 $\vx$ 支撑集中的第 $k$ 个超球面。
    存在一个标准正交基 $\vU_{k} \in \R^{D \times d_k}$ 使得 $\Span(\cS_k) = \Span(\vU_k)$。
    支撑集 $\cS$ 的以下正交分解将在整个证明中被反复使用。
    我们有
    \begin{align}
        \cS_{\delta}
        &= \set{\vxi \in \R^D \given \exists k \in [K] \::\: \dist(\vxi,
        \cS_k) \leq \delta}
        \\
        &= \bigcup_{k \in [K]} \set{\vxi \in \R^D \given \dist(\vxi,
        \cS_k) \leq \delta}.
    \end{align}
    通过正交投影，对于任意 $k \in [K]$，任何 $\vxi \in \R^D$ 都可以写成 $\vxi = \vxi^{\|} + \vxi^{\perp}$，其中 $\vxi^{\|} \in \Span(\cS_k)$ 且 $\ip{\vxi^{\|}}{\vxi^\perp} = 0$。
    那么对于任意 $\vxi' \in \cS_k$，我们有
    \begin{align}
        \norm*{\vxi - \vxi'}_2^2
        =
        \norm*{\vxi^{\|} + \vxi^{\perp} - \vxi'}_2^2
        &=
        \norm*{\vxi^{\|}}_2^2 + \norm*{\vxi^{\perp}}_2^2 + \norm*{\vxi'}_2^2
        - 2 \ip*{\vxi^{\|}}{\vxi'}
        \\
        &\geq
        \norm*{\vxi^{\|}}_2^2 + \norm*{\vxi'}_2^2
        - 2 \ip*{\vxi^{\|}}{\vxi'}
        \\
        &=
        \norm*{\vxi^{\|} - \vxi'}_2^2.
    \end{align}
    此外，已知对于任意非零的 $\vxi^{\|} \in \Span(\cS_k)$，
    \begin{equation}
        \inf_{\vxi' \in \cS_k}\,
        \norm*{\vxi^{\|} - \vxi'}_2^2
        =
        \norm*{\vxi^{\|} - \frac{\vxi^{\|}}{\norm{\vxi^{\|}}_2}}_2^2.
    \end{equation}
    如果 $\vxi^{\|}$ 为零，很明显对于每个 $\vxi' \in \cS_k$，上述距离等于 $1$。
    因此，如果我们定义一个投影映射 $\pi_{\cS_k}(\vxi)$ 如下
    \begin{equation}
        \pi_{S_k}(\vxi) = \frac{\vU_k\vU_k^\top \vxi}{\norm{\vU_k^\top \vxi}_2}
    \end{equation}
    对于任何满足 $\vU_k^\top \vxi \neq \mathbf{0}$ 的 $\vxi \in \R^D$，那么 $\pi_{\cS_k}(\vxi) = \argmin_{\vxi' \in \cS_k}\norm*{\vxi' - \vxi}_2$。
    我们选择 $0 < \delta < 1$，使得增厚集 $\cS_{\delta}$ 不包含任何使得投影映射 $\pi_{\cS_k}$ 未良定义的点 $\vxi \in \R^D$。
    因此增厚集 $\cS_{\delta}$ 满足
    \begin{align}
        \cS_{\delta}
        &= \bigcup_{k \in [K]} \set*{\vxi \in \R^D \given
            \norm*{\vxi - \frac{\vU_k\vU_k^\top
            \vxi}{\norm{\vU_k^\top \vxi}_2}}_2
        \leq \delta}.
    \end{align}
    这些距离可以用正交分解重写为
    \begin{align}
        \norm*{\vxi - \frac{\vU_k\vU_k^\top \vxi}{\norm{\vU_k^\top \vxi}_2}}_2^2
        &=
        \norm*{\vxi}_2^2 - 2 \norm{\vU_k^\top \vxi}_2 + 1
        \\
        &=
        \norm*{\vxi^{\|}}_2^2
        + \norm*{\vxi^{\perp}}_2^2
        - 2 \norm{\vU_k^\top \vxi^{\|}}_2 + 1
        \\
        &=
        \norm*{\vxi^{\perp}}_2^2
        + \left( \norm*{\vxi^{\|}}_2 - 1 \right)^2.
        \label{eq:distance-to-hypersphere}
    \end{align}
    接下来我们将证明，每一个这样的 $\vxi \in \cS_{\delta}$ 都可以唯一地关联到混合分布中单个子空间上的投影，这将允许我们定义一个到 $\cS$ 上的相应投影。
    给定一个 $\vxi \in \cS_{\delta}$，由上文可知，我们可以找到一个子空间 $\vU_k$，使得正交分解 $\vxi = \vxi^{\|}_k + \vxi^{\perp}_k$ 满足
    \begin{equation}
        \norm*{\vxi^{\perp}_k}_2^2
        + \left( \norm*{\vxi^{\|}_k}_2 - 1 \right)^2
        \leq
        \delta^2.
    \end{equation}
    考虑对于某个 $j \neq k$ 的分解 $\vxi = \vxi_j^{\|} + \vxi_j^{\perp}$。我们有
    \begin{align}
        \norm*{\vxi_j^{\|}}_2
        =
        \norm*{
            \vU_j \vU_j^\top \vxi
        }_2
        &=
        \norm*{
            \vU_j \vU_j^\top( \vU_k \vU_k^\top \vxi + (\vI - \vU_k \vU_k^\top)
            \vxi )
        }_2
        \\
        &=
        \norm*{
            \vU_j \vU_j^\top (\vI - \vU_k \vU_k^\top) \vxi
        }_2
        \\
        &\leq
        \norm*{
            (\vI - \vU_k \vU_k^\top) \vxi
        }_2
        =
        \norm*{
            \vxi_k^{\perp}
        }_2
        \leq \delta,
    \end{align}
    其中第二行使用了子空间 $\vU_k$ 的正交性假设，第三行使用了正交投影是非扩张的这一事实。
    因此，第 $j$ 个距离满足
    \begin{equation}
        \norm*{\vxi^{\perp}_j}_2^2
        + \left( \norm*{\vxi^{\|}_j}_2 - 1 \right)^2
        \geq
        (1 - \delta)^2.
    \end{equation}
    这意味着如果 $0 < \delta < 1/2$，每个 $\vxi \in \cS_{\delta}$ 在混合分布中都有一个唯一的最近子空间。
    因此，在此条件下，以下映射 $\pi_{\cS} : \cS_{\delta} \to \cS$ 是良定义的：
    \begin{equation}
        \pi_{\cS}(\vxi)
        =
        \pi_{\cS_{k_\star}}(\vxi),
        \enspace\text{where}\enspace
        k_{\star} = \argmin_{k \in [K]}\,
        \dist(\vxi, \cS_{k}).
    \end{equation}

    现在，我们为 $\vx_{\delta}$ 定义一个编码：
    \begin{equation}
        \hat{\vx}_{\delta} = \mathrm{q}( \pi_{\cS}(\vx_{\delta}) ).
    \end{equation}
    显然，这对应于 $\vx_{\delta}$ 的一个率 $R$ 编码，因为它使用了 $\vx$ 的率 $R$ 编码的编码-解码映射。我们必须证明它实现了较小的失真。
    我们计算
    \begin{align}
        \bE\left[ \norm*{ \vx_{\delta} - \hat{\vx}_{\delta} }_2^2 \right]
        &=
        \bE\left[ \norm*{ \vx_{\delta} -
        \mathrm{q}(\pi_{\cS}(\vx_{\delta})) }_2^2 \right]
        \\
        &\leq
        \left(
            \bE\left[ \norm*{ \vx_{\delta} - \pi_{\cS}(\vx_{\delta}) }_2^2
            \right]^{1/2}
            + \bE\left[ \norm*{ \pi_{\cS}(\vx_{\delta})
            - \mathrm{q}(\pi_{\cS}(\vx_{\delta})) }_2^2 \right]^{1/2}
        \right)^2,\label{eq:distortion-decomposition}
    \end{align}
    其中不等式使用了闵可夫斯基不等式（Minkowski inequality）。
    现在，根据 \Cref{def:thickening-set,def:thickening-rv-uos}，我们确定性地有
    \begin{equation}
        \norm*{ \vx_{\delta} - \pi_{\cS}(\vx_{\delta}) }_2^2
        \leq \delta^2,
    \end{equation}
    因此期望也满足这个估计。
    对于第二项，只需将随机变量 $\pi_{\cS}(\vx_{\delta})$ 的密度刻画为足够接近 $\vx$ 的密度即可——正如 \Cref{assumption:union-of-spheres} 所暗示的，$\vx$ 的密度是每个子球面 $\cS_k$ 上均匀分布的混合。
    根据上述论证，每个点 $\vxi \in \cS_{\delta}$ 都可以关联到唯一一个子空间 $\vU_k$，这意味着 $\cS_{\delta}$ 定义中的混合分量（回想 \Cref{def:thickening-rv-uos}）不重叠。因此，可以通过研究 $\pi_{\cS_k}$ 对条件随机变量 $\vx_{\delta}$（条件是其从 $\cS_{k, \delta}$ 中抽取）的影响来刻画密度 $\pi_{\cS}(\vx_{\delta})$。将此测度记为 $\mu_{k, \delta}$。
    我们断言，该测度在 $\pi_{\cS_k}$ 下的前推（pushforward）在 $\cS_k$ 上是均匀的。
    为了看出这一点成立，我们回顾 \Cref{eq:distance-to-hypersphere}，它给出了如下特征
    \begin{equation}
        \cS_{k, \delta} = \set*{
            \vxi^{\|} + \vxi^{\perp}
            \given
            \vxi^{\|} \in \Span(\vU_k),
            \vxi^{\perp} \in \Span(\vU_k)^\perp,
            \norm*{\vxi^{\perp}}_2^2
            + \left( \norm*{\vxi^{\|}}_2 - 1 \right)^2
            \leq
            \delta
        }.
    \end{equation}
    所讨论的条件分布在该集合上是均匀的；我们需要证明将投影 $\pi_{\cS_k}$ 应用于该条件随机变量会产生一个在 $\cS_k$ 上均匀的随机变量。
    关于这些坐标，我们已经看到 $\pi_{\cS_k}(\vxi^\| + \vxi^\perp) = \vxi^\| / \norm{\vxi^{\|}}_2$。
    因此，对于任意 $\vxi \in \cS_{k}$，$\vxi$ 在 $\pi_{\cS_k}$ 下于 $\cS_{k, \delta}$ 中的原像为
    \begin{equation}
        \pi_{\cS_k}^{-1}(\vxi) = \set*{
            r\vxi + \vxi^\perp
            \given
            r > 0,
            \vxi^{\perp} \in \Span(\vU_k)^\perp,
            \norm*{\vxi^{\perp}}_2^2
            + \left( r - 1 \right)^2
            \leq
            \delta
        }.
        \label{eq:fiber-of-uos}
    \end{equation}
    为了证明 $(\pi_{\cS_k})_{\sharp} \mu_{k, \delta}$ 是均匀的，我们需要以一种方式分解 $\cS_{k, \delta}$ 上均匀密度的积分，使得清楚地看出每个纤维 $\pi_{\cS_k}^{-1}(\vxi)$（对于每个 $\vxi \in \cS_k$）对积分的“贡献”是相等的。\footnote{更严格地说，这对应于将 $\cS_{k, \delta}$ 上的均匀密度分解为对应于 $\vxi \in \cS_{k}$ 的正则条件密度，并证明相应的关于 $\vxi$ 的密度是均匀的。证明清楚地表明这是正确的。}
    根据 \Cref{def:thickening-rv-uos}，我们有
    \begin{equation}
        \volume(\cS_{k, \delta})
        = \iint_{\Span(\vU_k) \times \Span(\vU_k)^\perp} \mathbf{1}_{
            \norm*{\vxi^{\perp}}_2^2
            + \left( \norm*{\vxi^{\|}}_2 - 1 \right)^2
            \leq
            \delta
        }
        \odif \vxi^{\|} \odif \vxi^{\perp}.
    \end{equation}
    特别地，关于正交坐标的积分可以分解。
    令 $\odif \vtheta^{d}$ 表示 $\R^d$ 中半径为 $1$ 的球面上的均匀（Haar）测度。将 $\vxi^{\|}$ 积分转换为极坐标，我们有
    \begin{equation}
        \volume(\cS_{k, \delta})
        = \int_{[0, \infty)} \int_{\bS^{d_k-1}} \int_{\Span(\vU_k)^\perp}
        r^{d_k-1}
        \mathbf{1}_{
            \norm*{\vxi^{\perp}}_2^2
            + \left( r - 1 \right)^2
            \leq
            \delta
        }
        \odif r \odif \vtheta^{d_k} \odif \vxi^{\perp}.
    \end{equation}
    与纤维表示 \eqref{eq:fiber-of-uos} 相比，我们看到我们需要对前面积分的 $r$ 和 $\vxi^\perp$ 分量进行“积分消去”（integrate out），以验证前推是均匀的。
    但这显而易见，因为前面的表达式表明该积分的值独立于 $\vxi^\|$——或者在上下文中等价地，独立于球面分量 $\vtheta^{d_k}$ 的值。

    因此，由上述论证可知 $\pi_{\cS}(\vx_{\delta})$ 是均匀的。
    因为对 $\delta$ 的假设意味着 $\vx_{\delta}$ 分布中的混合分量不重叠，所以混合权重 $\pi_k$ 在像 $\pi_{\cS}(\vx_{\delta})$ 中也得以保留，特别地，$\pi_{\cS}(\vx_{\delta})$ 的分布等于 $\vx$ 的分布。
    因此，\Cref{eq:distortion-decomposition} 中的第二项满足
    \begin{equation}
        \bE\left[ \norm*{ \pi_{\cS}(\vx_{\delta}) -
        \mathrm{q}(\pi_{\cS}(\vx_{\delta})) }_2^2 \right]
        =
        \bE\left[ \norm*{ \vx - \mathrm{q}(\vx) }_2^2 \right]
        \leq
        \epsilon^2,
    \end{equation}
    因为 $\mathrm{q}$ 是 $\vx$ 的失真为 $\epsilon$ 的编码。

    因此我们已经证明，假设的 $\vx$ 的率 $R$、（期望平方）失真 $\epsilon^2$ 编码产生了一个 $\vx_{\delta}$ 的率 $R$、（期望平方）失真 $\delta + \epsilon$ 编码。
    这确立了
    \begin{equation}
        R_{\delta + \epsilon}(\vx_{\delta})
        \leq
        R_{\epsilon}(\vx),
    \end{equation}
    证毕。

\end{proof}

\end{document}